\chapter{The threshold inputs for general pure models}\label{app:general-selection}

In this section, we prove \cref{prop:threshold-inputs} for every $p\ge3$.
For $p\in\{3,4\}$, \cref{sec:tree-ceilings} proves both parts of that
proposition by a deterministic inequality for tree vectors. The argument
given here does not use the tree ceilings and works for every pure model.
Part~\ref{item:threshold-ceiling} of \cref{prop:threshold-inputs} is the
input to the proof of marginality in \cref{sel:positive-marginality}, and
part~\ref{item:threshold-forgetting} is the input to the proof of forgetting
in \cref{sel:fixed-wait-forgetting}. The finite forks of \cref{sec:general-forks}, the overlap penalty of
\cref{sel:absolute-field-cost}, and the saturation identity
\eqref{sel:fork-saturation} are used again in \cref{resp:appendix-p}, where
we extend the localization of the response to every $p\ge3$.

We use the landscape bounds of \cref{lem:landscape-moments}, the fixed-horizon bridge of \cref{prop:hard-bridge}, and the Brownian form of the
averaged response in \cref{lem:brown-force-response}. From
\cref{sec:ck-structure}, we use the energy identity of \cref{prop:energy-bridge} and the bound
\eqref{eq:finite-lag-flow-moments} on the moments of the response, proved there. From
\cref{sec:entropy-clusters}, we use the entropy budget of
\cref{lem:entropy-budget}, the integration by parts formula
\eqref{eq:gibbs-current-integration-by-parts}, the Fisher majorant of
\cref{thm:fisher-clock}, and the stationary clusters of
\cref{prop:stationary-clusters}. We also use the formulas
\eqref{sel:prefix-norm} and \eqref{sel:prefix-energy} for $d_L$ and $h_L$.
The proof of part~\ref{item:threshold-forgetting} uses, in addition, the
identities \eqref{sel:limiting-center-energy} and \eqref{sel:center-amgm}
and the correlation modulus \eqref{eq:global-modulus}.

\paragraph{Ideas of the proof.}
We obtain the energy ceiling in part~\ref{item:threshold-ceiling} from the
algorithmic threshold of Huang and Sellke. At a late time $T$, we continue
the trajectory by $K$ independent Brownian futures of length $L$, which we
call a fork, and average their endpoints. The current identity of
\cref{sec:entropy-clusters} shows that two branches of the fork have
overlap close to $c(a+d)$, where $a$ and $d$ are their lengths. We compute
the energy of the average exactly by Gaussian integration by parts.
Subtracting the energy of a single branch removes the common past in the
dimension limit. Using the hard--soft approximation of Sellke, we approximate the average by
a function of the disorder that is Lipschitz with a constant independent of
the dimension. The algorithmic threshold then bounds its energy. Letting
$K\to\infty$, we obtain part~\ref{item:threshold-ceiling}.

For part~\ref{item:threshold-forgetting}, we add a deterministic field to
the threshold. By convex duality, the field version of the threshold
penalizes the overlap of a near-threshold algorithm with any direction
that is independent of the disorder. However, at positive times, an earlier state $X_s$ depends on the disorder. We therefore disclose the disorder in stages,
along a finite sequence of gradient queries, and apply the penalty to one
new direction at each stage. An Euler--Maruyama approximation places $X_s$
close to the span
of finitely many such queries. Under marginality, the normalized average of
a fork saturates the threshold, and the penalty forces its overlap with
$X_s$ to vanish. Up to normalization, this overlap is $C(t,s)$.

\paragraph{Attribution.}
The algorithmic threshold is \cite[Theorem~1 and Corollary~1.8]{hs2023};
for even $p$ it was proved earlier in \cite{hs2021}. In
\cite[Section~3]{sellke2024}, Sellke combined it with a Lipschitz
approximation of Langevin dynamics on a fixed time interval to show that
the energy of the dynamics stays below the threshold. We follow that
argument in \cref{sec:general-forks}. Trajectories with a common past and
independent future noise are a standard construction in the physics of
aging. Cugliandolo and Dean considered replicas that share their noise up
to a waiting time and receive independent noise afterward in
\cite[arXiv version~1, Section~II]{cugliandolodean1995}. Barrat, Burioni, and M\'ezard classified aging through the overlaps of such copies in
\cite[Sections~2 and~4]{bbm1996}. Franz and Virasoro
proposed averaging cloned futures to estimate a quasi-state magnetization in
\cite[arXiv version~3, Section~III]{franzvirasoro2000}, and noise-averaged
magnetizations are central to Biroli's dynamical TAP approach in
\cite[arXiv version~2, Section~III.B]{biroli1999}. Here we prove the overlap and energy identities of the forks from the
current identity of \cref{sec:entropy-clusters} and from Gaussian
integration by parts. The
disclosure argument of \cref{sec:general-disclosure} is a Gaussian
conditioning argument; we credit its antecedents there.

\paragraph{Order of limits.}
Every horizon, number of branches, field strength, approximation accuracy,
and Lipschitz constant below is fixed while the dimension tends to
infinity, and we take the long-time and outer limits only afterward.
We take dimension limits along the subsequence $(N_j)$ of
\cref{thm:fisher-clock}, and we often write $N\to\infty$ for $N=N_j$,
$j\to\infty$. The parameters and their order are as follows:
\begin{equation}\label{sel:aging-limit-order}
 N_j\to\infty,\qquad n\to\infty,\qquad \varepsilon\downarrow0,\qquad
 K\to\infty,\qquad L\to\infty,\qquad M\to\infty.
\end{equation}
Here $n$ indexes a sequence of late times ($T_n$ in
\cref{sec:general-forks} and $t_n$ in \cref{sec:general-transfer}),
$\varepsilon$ is the accuracy of the Lipschitz approximation in
\cref{sel:finite-forks}, $K$ is the number of branches of a fork, $L$ is
their length, and $M$ is the number of steps of the Euler--Maruyama scheme of
the proof of \cref{sel:general-fixed-wait}. Each parameter is fixed while the limits to its
left are taken. Part~\ref{item:threshold-ceiling} uses only the first four
limits, at a fixed $L$; part~\ref{item:threshold-forgetting} uses all six.
For the forgetting proof, it suffices to take the approximation accuracies
$\varepsilon=1/(m+1)$, $m\in\mathbb N$, so $\varepsilon\downarrow0$ there
means $m\to\infty$.
We remove the tolerance and the field strengths of the threshold theorem
inside the proof of \cref{sel:absolute-field-cost}, after the dimension
limit at fixed values of all other parameters. Constants of the Lipschitz
maps may depend on every fixed parameter, and they are used only before the
limits in which that parameter varies.

\section{The algorithmic threshold and its overlap penalty}
\label{sec:general-field-threshold}

Throughout this section, $\norm{\cdot}$ is the Euclidean norm,
\[
 B_N=\{x\in\R^N:\norm x\le\sqrt N\},
\]
and $\Pi_{B_N}$ is the metric projection onto $B_N$, which is $1$-Lipschitz.
The coefficient tensor $J$ of \eqref{eq:hamiltonian} lies in the
$N^p$-dimensional space $(\R^N)^{\otimes p}$, with its Euclidean norm, and
has the standard Gaussian law there. More explicitly, we keep the ordered
coefficients and use the Hilbert norm
\[
 \norm J_{\mathrm{coef}}^2
 =\sum_{i_1,\dots,i_p=1}^N J_{i_1\dots i_p}^2 .
\]
This norm is not rescaled with the dimension. When the dependence on $J$ matters, we
write $H_{N,J}$ for $H_N$. An algorithm is a measurable map
$\mathcal A_N=\mathcal A_N(J,\zeta)$ with values in $B_N$, where the seed
$\zeta$ is a random variable independent of $J$ (possibly absent). For
$0\le L_{\mathrm{Lip}}<\infty$, we say that $\mathcal A_N$ is $L_{\mathrm{Lip}}$-Lipschitz if, for
almost every value of $\zeta$,
\[
 \norm{\mathcal A_N(J,\zeta)-\mathcal A_N(J',\zeta)}\le L_{\mathrm{Lip}}\norm{J-J'}_{\mathrm{coef}}
 \qquad\text{for all }J,J'\in(\R^N)^{\otimes p}.
\]
The constant $L_{\mathrm{Lip}}$ is deterministic and does not depend on $N$ or on the
seed.

Recall the tensor norm $Q_N$ of \eqref{eq:tensor-norm} and the constants
$A_p,c_p>0$ of \cref{lem:landscape-moments}, for which
$\Prob(Q_N>A_p)\le e^{-c_pN}$. We write
\[
 Q_*=\sup_N(\E Q_N^2)^{1/2},
\]
which is finite by \cref{lem:landscape-moments}. By the definition of
$Q_N$, a contraction of $J$ against vectors $u_1,\dots,u_p$ is at most
$\sqrt N\,Q_N\norm{u_1}\cdots\norm{u_p}$ in absolute value. A derivative of
order $k\le2$ of $H_N$ at $x$, applied to $k$ vectors, is
$N^{-(p-1)/2}$ times a sum of $p(p-1)\cdots(p-k+1)$ such contractions in
which $p-k$ slots equal $x$. For $R\ge1$ and $\norm x\le R\sqrt N$, this
implies that
\begin{equation}\label{eq:general-ball-derivatives}
\begin{gathered}
 \frac{|H_N(x)|}N\le R^pQ_N,\qquad
 \norm{\nabla H_N(x)}\le pR^{p-1}Q_N\sqrt N,\\
 \norm{\nabla^2H_N(x)}_\op\le p(p-1)R^{p-2}Q_N.
\end{gathered}
\end{equation}
For $R=1$, the first two bounds read
\begin{equation}\label{sel:tensor-ui}
 \sup_{x\in B_N}\frac{|H_N(x)|}N\le Q_N,\qquad
 \sup_{x\in B_N}\norm{\nabla H_N(x)}\le pQ_N\sqrt N.
\end{equation}
In particular, the normalized energy of any algorithm is bounded by $Q_N$,
whose moments are bounded uniformly in $N$. Now let $x$ and $y$ be random
vectors of norm at most $R\sqrt N$. The segment between them lies in the
same ball. Using the mean value theorem, \eqref{eq:general-ball-derivatives},
and the Cauchy--Schwarz inequality, we find
\begin{equation}\label{eq:general-energy-transfer}
 \frac{\E|H_N(x)-H_N(y)|}N
 \le pR^{p-1}Q_*\left(\E\frac{\norm{x-y}^2}N\right)^{1/2}.
\end{equation}

For $\kappa\ge0$, let $q_1(\kappa)=(\kappa^2/(p(p-2)))^{1/(p-1)}$, and
define
\begin{equation}\label{eq:alg-field}
 \ALG_p(\kappa)=
 \begin{cases}
 \Eth+\dfrac{p-2}{p}\sqrt{p(p-1)}\,q_1(\kappa)^{p/2},
   &\kappa^2\le p(p-2),\\[3mm]
 \sqrt{p+\kappa^2},&\kappa^2\ge p(p-2).
 \end{cases}
\end{equation}
In terms of $\kappa$, the first line equals
$\Eth+\gamma_p\kappa^{p/(p-1)}$ with
$\gamma_p=((p-2)/p)\sqrt{p(p-1)}\,[p(p-2)]^{-p/(2(p-1))}$.

\begin{theorem}[External-field Lipschitz threshold of Huang and Sellke]
\label{sel:published-field-threshold}
Fix an integer $p\ge3$, $0\le L_{\mathrm{Lip}}<\infty$, $\kappa\ge0$, and $\eta>0$. There exist $c>0$
and $N_0$, depending only on $p$, $L_{\mathrm{Lip}}$, $\kappa$, and $\eta$, such that
the following holds for every $N\ge N_0$. Let $w\in\R^N$ be deterministic
with $\norm w=\sqrt N$, and let $\mathcal A_N:(\R^N)^{\otimes p}\to B_N$ be
$L_{\mathrm{Lip}}$-Lipschitz. Then
\begin{equation}\label{sel:field-threshold}
 \Prob\left(\frac{H_N(\mathcal A_N(J))+\kappa\ip{w}{\mathcal A_N(J)}}N
 \ge\ALG_p(\kappa)+\eta\right)\le e^{-cN}.
\end{equation}
\end{theorem}

\Cref{sel:published-field-threshold} is the case of one species and
$\xi(q)=q^p$ of \cite[Theorem~1 and Corollary~1.8]{hs2023}. We now check
that our setting satisfies the hypotheses there.
In \cite{hs2023}, the Hamiltonian is $\ip{h}{\sigma}$ plus a Gaussian part
on the ball of radius $\sqrt N$, the field is $h$ times the all-ones vector
$\mathbf 1$ with $h\ge0$, and an algorithm is a function of the disorder
coefficients, concatenated into one vector with the Euclidean distance.
Theorem~1 there states that, for fixed $\tau,\eta>0$ and $N$ large, every
$\tau$-Lipschitz algorithm reaches energy at least $\ALG+\eta$ with
probability at most $e^{-cN}$, with $c$ and the threshold for $N$ depending
on $\tau$, $\eta$, the covariance, and the field only. For one species with
$\xi(q)=q^p$, the Gaussian part is exactly $H_N$ of \eqref{eq:hamiltonian},
whose ordered coefficients are independent standard Gaussians; we take
$h=\kappa$. Our algorithms are functions of $J$ alone, the field being
deterministic. Corollary~1.8 of \cite{hs2023}, which covers odd and even
$p$, states that $\ALG=(\xi'(1)+\kappa^2)^{1/2}$ if
$\xi'(1)+\kappa^2\ge\xi''(1)$, and otherwise
\[
 \ALG=q_1\,\xi''(q_1)^{1/2}+\int_{q_1}^1\xi''(q)^{1/2}\,\dd q,
\]
where $q_1\in[0,1)$ solves $\xi'(q_1)+\kappa^2=q_1\xi''(q_1)$. For
$\xi(q)=q^p$, the first condition reads $\kappa^2\ge p(p-2)$, the equation
for $q_1$ reads $\kappa^2=p(p-2)q_1^{p-1}$, and
\[
\begin{aligned}
 q_1\xi''(q_1)^{1/2}+\int_{q_1}^1\xi''(q)^{1/2}\,\dd q
 &=\sqrt{p(p-1)}\left(q_1^{p/2}+\frac2p\big(1-q_1^{p/2}\big)\right)\\
 &=\Eth+\frac{p-2}p\sqrt{p(p-1)}\,q_1^{p/2},
\end{aligned}
\]
where we used $2\sqrt{p(p-1)}/p=\Eth$. This is \eqref{eq:alg-field}. At
$\kappa^2=p(p-2)$, we have $q_1=1$ and both lines of \eqref{eq:alg-field}
equal $\sqrt{p(p-1)}$.

It remains to replace $\mathbf 1$ by an arbitrary direction $w$. Let $O$ be
an orthogonal matrix such that $O\mathbf 1=w$. The map $J\mapsto O^{\otimes p}J$ is
a linear isometry of $(\R^N)^{\otimes p}$ that preserves the standard
Gaussian law, and $H_{N,O^{\otimes p}J}(Ox)=H_{N,J}(x)$ for all $x$. Given
$\mathcal A_N$, the map $\mathcal A'_N(J)=O^{\mathsf T}\mathcal A_N(O^{\otimes p}J)$
is $L_{\mathrm{Lip}}$-Lipschitz with values in $B_N$, and, with $J'=O^{\otimes p}J$,
\[
 H_{N,J}(\mathcal A'_N(J))+\kappa\ip{\mathbf 1}{\mathcal A'_N(J)}
 =H_{N,J'}(\mathcal A_N(J'))+\kappa\ip{w}{\mathcal A_N(J')}.
\]
The tensor $J'$ has the law of $J$, and \cite[Theorem~1]{hs2023} applied to
$\mathcal A'_N$ gives \eqref{sel:field-threshold}, with constants that do
not depend on $w$ or on $\mathcal A_N$. For $\kappa=0$, we have
$\ALG_p(0)=\Eth$, and \cref{sel:published-field-threshold} is the
pure-model threshold, recorded in this form in
\cite[Proposition~1.9]{sellke2024}.

In this argument, we may always use the larger positive bound
$L_{\mathrm{Lip}}+1$. This includes
$L_{\mathrm{Lip}}=0$ and gives the same uniform choice of $c$ and $N_0$
for every algorithm with the prescribed original bound. We need no separate
measurability assumption on a deterministic algorithm, since Lipschitz maps
are continuous.

\paragraph{Constructive finite-tree qualifications.}
The branching proof of the threshold uses finite trees whose depth and
branching degrees are fixed before the dimension tends to infinity. Two
auxiliary statements in \cite[Appendix~A]{hs2023} require qualifications
in a constructive version. In Lemma~A.8, if the absolute values of the leaf
energies are bounded by $C>0$ and the depth is $D\ge1$, the retained branching
degree $\lfloor K\varepsilon/(3CD)\rfloor$ is obtained in the range
$0<\varepsilon\le CD$, and $K$ is then taken large enough to make this
degree positive. This is the range used when $\varepsilon\downarrow0$;
an unrestricted error can make the displayed degree exceed $K$.

For the base-preserving fork correction of Lemma~A.9, we impose $k<N$
when $k$ orthogonal increments are retained. After dividing
the vertices by $\sqrt N$, suppose that the base $x$ has squared norm
$q\ge q_0>0$, the target increment has squared norm $q'\ge0$, and the
approximate increments $y_1,\dots,y_k$ satisfy, for $\varepsilon\ge0$,
\[
 |\norm{y_i}^2-q'|\le\varepsilon,\qquad
 |\ip x{y_i}|\le\varepsilon,\qquad
 |\ip{y_i}{y_j}|\le\varepsilon\quad(i\ne j).
\]
Under the sufficient small-error conditions
$\sqrt\varepsilon\le q_0$ and $4(k+1)\sqrt\varepsilon\le1$, the base can
be kept exactly and the increments corrected to mutually orthogonal
vectors $z_i\perp x$ with $\norm{z_i}^2=q'$ and
$\norm{z_i-y_i}\le e_k(\varepsilon)$, where
\[
 e_k(\varepsilon)
 =(4\sqrt{k+1}+1)\sqrt\varepsilon
   +2\sqrt{\sqrt\varepsilon}.
\]
This bound permits $q'=0$ and tends to zero for fixed $k$ as
$\varepsilon\downarrow0$. It supplies a vanishing modulus, independent of the dimension, in place of
the piecewise numerical modulus printed in that lemma.

We handle a schedule with zero initial overlap before this correction.
For fixed $\delta\in(0,1)$, we scale the normalized vertices, which have norm
at most one, by $\sqrt{1-\delta}$ and add a common orthogonal vector of norm
$\sqrt\delta$. The new schedule of target overlaps is
$q_i^{(\delta)}=(1-\delta)q_i+\delta$, so its root is at least $\delta$;
the terminal vectors change by at most $2\sqrt\delta$. We choose the dimension larger than the number of vertices of the retained
finite tree, as well as its branching degree, so that the common orthogonal
direction and all fork corrections exist. For fixed depth and retained
branching, the accumulated correction tends to zero, and $\delta$ can be
removed afterward. Repeated overlap levels are allowed: their target
increments have norm zero.

In the algorithmic application, we can retain a fixed deterministic set of children, because the success event imposes the desired energy bound
at every original terminal leaf. Internal vertices are first obtained
from averages of the original terminal outputs. Restricting the independent Gaussian increments of the tree to this fixed
set preserves their joint law. We therefore need no distributional assertion
about an adaptively selected subtree. The depths, branching degrees, error
tolerances, and dimension cutoffs are chosen before the algorithm and the
deterministic field direction. The calibrated overlap schedule may depend on the algorithm,
since the finite-tree ceiling is uniform over all admissible schedules.
This is the uniformity asserted in \cref{sel:published-field-threshold}.

The field version of the threshold penalizes the overlap of an algorithm
with a direction independent of $J$. To state the penalty, recall
$\qd=(p-2)/(p-1)$ and $\bd^2=(p-1)^{p-1}/(p(p-2)^{p-2})$. The constant $\bd$
is the dynamical transition point, and it enters here only as the
coefficient of the penalty. The inverse temperature $\beta$ of the dynamics
plays no role in \cref{sel:published-field-threshold,sel:absolute-field-cost}.

\begin{lemma}[Absolute-overlap energy cost]\label{sel:absolute-field-cost}
Fix an integer $p\ge3$ and $0\le L_{\mathrm{Lip}}<\infty$. There exist
finite real numbers $\phi_N=\phi_N(p,L_{\mathrm{Lip}})\ge0$, for every $N\ge1$,
with $\phi_N\to0$ as $N\to\infty$
such that the following holds. Let $\zeta$ be a random variable independent of $J$, let
$w=w(\zeta)\in\R^N$ satisfy $\norm w=\sqrt N$, let $\mathcal A_N$ be an
$L_{\mathrm{Lip}}$-Lipschitz algorithm with seed $\zeta$, and set
$a_N=\ip{w}{\mathcal A_N}/N$. Then, almost surely,
\begin{equation}\label{eq:conditional-field-cost}
 \E\left[\frac{H_N(\mathcal A_N)}N+\bd|a_N|^p\,\middle|\,\zeta\right]
 \le\Eth+\phi_N.
\end{equation}
In particular,
\begin{equation}\label{sel:absolute-field-cost-equation}
 \limsup_{N\to\infty}
 \left\{\E\frac{H_N(\mathcal A_N)}N+\bd\,\E|a_N|^p\right\}\le\Eth.
\end{equation}
\end{lemma}

The lemma makes no assumption on the concentration of $a_N$, and the seed
may be absent, in which case $w$ is deterministic.

\begin{proof}
The proof has three steps. We first compute the best ceiling that the field
thresholds give for a given overlap, then we use finitely many fields at
once, and finally we refine the set of fields after the dimension limit.

\paragraph{Step 1: the conjugate ceiling.}
For $0\le a\le1$, we define
\[
 \mathcal E_p(a)=\inf_{\kappa\ge0}\{\ALG_p(\kappa)-\kappa a\}.
\]
We claim that
\begin{equation}\label{eq:conjugate-ceiling}
 \mathcal E_p(a)=
 \begin{cases}
 \Eth-\bd a^p,&0\le a\le\sqrt{\qd},\\
 \sqrt{p(1-a^2)},&\sqrt{\qd}\le a\le1,
 \end{cases}
 \qquad\text{and}\qquad
 \Eth-\mathcal E_p(a)\ge\bd a^p\quad(0\le a\le1).
\end{equation}
The function $\ALG_p$ is convex on $[0,\infty)$, since $\kappa^{p/(p-1)}$ and
$\sqrt{p+\kappa^2}$ are convex. It is also continuously differentiable. Indeed, let $\kappa_*=\sqrt{p(p-2)}$. On $[0,\kappa_*]$, its derivative
is $p(p-1)^{-1}\gamma_p\kappa^{1/(p-1)}$, which increases from $0$ to
\[
 \frac p{p-1}\gamma_p\kappa_*^{1/(p-1)}
 =\frac{p-2}{p-1}\sqrt{\frac{p(p-1)}{p(p-2)}}=\sqrt{\qd}.
\]
On $[\kappa_*,\infty)$, its derivative is $\kappa/\sqrt{p+\kappa^2}$,
which increases from $\sqrt{\qd}$ to $1$. By convexity, for $0\le a<1$ the
infimum defining $\mathcal E_p(a)$ is attained where $\ALG_p'(\kappa)=a$.

If $a\le\sqrt{\qd}$, this point is
$\kappa=((p-1)a/(p\gamma_p))^{p-1}$. At this point, we have $\gamma_p\kappa^{1/(p-1)}=(p-1)a/p$, and therefore
\[
 \mathcal E_p(a)=\Eth+\kappa\big(\gamma_p\kappa^{1/(p-1)}-a\big)
 =\Eth-\frac{\kappa a}p
 =\Eth-\frac1p\left(\frac{p-1}{p\gamma_p}\right)^{p-1}a^p.
\]
Inserting the value of $\gamma_p$, we have
\[
\begin{aligned}
 \left(\frac{p-1}{p\gamma_p}\right)^{p-1}
 &=\left(\frac{p-1}{p-2}\right)^{p-1}\big(p(p-1)\big)^{-(p-1)/2}
  \big(p(p-2)\big)^{p/2}\\
 &=p^{1/2}(p-1)^{(p-1)/2}(p-2)^{-(p-2)/2},
\end{aligned}
\]
and dividing by $p$, we find that the coefficient of $a^p$ is
$((p-1)^{p-1}/(p(p-2)^{p-2}))^{1/2}=\bd$.

If $\sqrt{\qd}\le a<1$, the point is $\kappa=\sqrt p\,a/\sqrt{1-a^2}$, and
\[
 \mathcal E_p(a)=\sqrt{\frac p{1-a^2}}-\frac{\sqrt p\,a^2}{\sqrt{1-a^2}}
 =\sqrt{p(1-a^2)}.
\]
At $a=1$, we have $\mathcal E_p(1)=0$, because the function
$\sqrt{p+\kappa^2}-\kappa$ decreases to $0$ as $\kappa\to\infty$. The two formulas agree at
$a=\sqrt{\qd}$, where both equal $\sqrt{p/(p-1)}$. This proves the first
part of \eqref{eq:conjugate-ceiling}, and shows that $\mathcal E_p$ is
continuous on $[0,1]$.

Let $D_p(a)=\Eth-\mathcal E_p(a)$. For $a\le\sqrt{\qd}$, we have
$D_p(a)=\bd a^p$. For $\sqrt{\qd}<a<1$, we have
$D_p'(a)=\sqrt p\,a/\sqrt{1-a^2}$. The function $a^{p-2}\sqrt{1-a^2}$ on
$[0,1]$ is maximal at $a^2=\qd$, where it equals
\[
 \qd^{(p-2)/2}(1-\qd)^{1/2}=(p-2)^{(p-2)/2}(p-1)^{-(p-1)/2}
 =\frac{1}{\sqrt p\,\bd}.
\]
Since $D_p'(a)/(p\bd a^{p-1})=(\sqrt p\,\bd\,a^{p-2}\sqrt{1-a^2})^{-1}$, we conclude that $D_p'(a)\ge p\bd a^{p-1}$ on $(\sqrt{\qd},1)$. Integrating this
inequality from $\sqrt{\qd}$, where $D_p(\sqrt{\qd})=\bd\qd^{p/2}$, and
using the continuity of $D_p$ at $a=1$, we find that $D_p(a)\ge\bd a^p$ on
$[\sqrt{\qd},1]$. This proves \eqref{eq:conjugate-ceiling}.

\paragraph{Step 2: finitely many fields.}
Fix a finite set $\mathcal G\subset[0,\infty)$ with $0\in\mathcal G$, and
$\eta>0$. We define
\[
 V_{\mathcal G}(a)=\min_{\kappa\in\mathcal G}\{\ALG_p(\kappa)-\kappa a\},
 \qquad 0\le a\le1.
\]
Note that $|V_{\mathcal G}|\le M_{\mathcal G}$, where
$M_{\mathcal G}=\Eth+\max\mathcal G$, because $0\in\mathcal G$ and
$\ALG_p\ge0$. We condition on $\zeta$. Then $w$ is
fixed, $J$ is standard Gaussian, and $J\mapsto\mathcal A_N(J,\zeta)$ is
$L_{\mathrm{Lip}}$-Lipschitz. For each $\kappa\in\mathcal G$, we use
\cref{sel:published-field-threshold} with the field strength $\kappa$ and the
directions $w$ and $-w$. Let $c(\mathcal G,\eta)>0$ and $N_0(\mathcal G,\eta)$
be the smallest constant and the largest threshold for $N$ among these
$|\mathcal G|$ applications. For $N\ge N_0(\mathcal G,\eta)$, the event
\[
 \mathfrak B=\bigcup_{\kappa\in\mathcal G}\left\{
 \frac{H_N(\mathcal A_N)}N+\kappa|a_N|\ge\ALG_p(\kappa)+\eta\right\}
\]
has conditional probability at most $2|\mathcal G|e^{-c(\mathcal G,\eta)N}$,
uniformly in $\zeta$. Off $\mathfrak B$, we have
$H_N(\mathcal A_N)/N<V_{\mathcal G}(|a_N|)+\eta$; note that $|a_N|\le1$,
since $\norm w=\sqrt N$ and $\mathcal A_N\in B_N$. On $\mathfrak B$, we use
$|H_N(\mathcal A_N)|/N\le Q_N$ from \eqref{sel:tensor-ui}. Recall that $J$ is
independent of $\zeta$, which gives $\E[Q_N^2\mid\zeta]=\E Q_N^2\le Q_*^2$.
Let
\[
 r_N(\mathcal G,\eta)=(Q_*+M_{\mathcal G})
 \big(2|\mathcal G|e^{-c(\mathcal G,\eta)N}\big)^{1/2}.
\]
Using the Cauchy--Schwarz inequality, we obtain
\begin{equation}\label{eq:finite-grid-cost}
 \E\left[\frac{H_N(\mathcal A_N)}N-V_{\mathcal G}(|a_N|)\,\middle|\,\zeta\right]
 \le\eta+r_N(\mathcal G,\eta)
\end{equation}
for $N\ge N_0(\mathcal G,\eta)$. For $1\le N<N_0(\mathcal G,\eta)$, we set
$r_N(\mathcal G,\eta)=Q_*+M_{\mathcal G}$. Then
\eqref{eq:finite-grid-cost} still holds, directly from
$|H_N(\mathcal A_N)|/N\le Q_N$ and $|V_{\mathcal G}|\le M_{\mathcal G}$.
The finite quantity $r_N(\mathcal G,\eta)$ depends only on $p$, $L_{\mathrm{Lip}}$, $\mathcal G$, and $\eta$, and tends to zero as $N\to\infty$.

\paragraph{Step 3: refining the fields.}
Let $\mathcal G_k=\{i2^{-k}:0\le i\le4^k\}$. These sets increase with $k$,
and their union is dense in $[0,\infty)$. The functions
$V_{\mathcal G_k}$ are continuous on $[0,1]$ and decrease in $k$. For each
$a\in[0,1]$, the map $\kappa\mapsto\ALG_p(\kappa)-\kappa a$ is continuous,
and its infimum over the dense set $\bigcup_k\mathcal G_k$ equals
$\mathcal E_p(a)$. The limit $\mathcal E_p$ is continuous on $[0,1]$. By
Dini's theorem, the quantity
$s_k=\sup_{0\le a\le1}\{V_{\mathcal G_k}(a)-\mathcal E_p(a)\}$ tends to zero.
By \eqref{eq:conjugate-ceiling} and \eqref{eq:finite-grid-cost} with
$\mathcal G=\mathcal G_k$ and $\eta=1/k$, we have
\[
\begin{aligned}
 \E\left[\frac{H_N(\mathcal A_N)}N+\bd|a_N|^p\,\middle|\,\zeta\right]
 &\le\Eth+\E\left[\frac{H_N(\mathcal A_N)}N-\mathcal E_p(|a_N|)\,\middle|\,\zeta\right]\\
 &\le\Eth+s_k+\frac1k+r_N(\mathcal G_k,1/k)
\end{aligned}
\]
for every $k\ge1$. We define
$\phi_N=\inf_{k\ge1}\{s_k+1/k+r_N(\mathcal G_k,1/k)\}$, which depends only on
$p$ and $L_{\mathrm{Lip}}$. This infimum is nonnegative and finite for every
$N\ge1$, since already the term with $k=1$ is finite. Then
\eqref{eq:conditional-field-cost} holds. Moreover,
$\limsup_N\phi_N\le s_k+1/k$ for every $k$, and in particular $\phi_N\to0$. Taking
expectations in \eqref{eq:conditional-field-cost} and letting $N\to\infty$
gives \eqref{sel:absolute-field-cost-equation}.
\end{proof}

\section{Finite forks and the exact energy subtraction}\label{sec:general-forks}

In this section, we fix $\beta>0$ and a stationary cluster $(c,r,\ell)$ of
the translated CK kernels, as in \cref{prop:stationary-clusters}, realized
along a sequence $T_n\to\infty$. Its plateau is $q$ and its energy is
$e_\infty=(\ell-1)/(p\beta)$. Recall from \eqref{eq:cluster-fdt-properties}
that $0\le c\le1$, that $r=-c'\ge0$, so that $c$ is nonincreasing, and that
$\int_0^\infty r(\tau)\,\dd\tau=1-q$. Recall also that $e_\beta(T_n+a)\to e_\infty$
for every $a\ge0$, by \eqref{eq:energy-identity} and the convergence of
the translated multipliers.

\paragraph{Entropy notation.}
We use the notation of \cref{sec:entropy-clusters}. For fixed disorder,
$\rho_t$ is the density of $X_t$ with respect to $\omega_N$ and $g_N$ is the
Gibbs density. We write
\[
 v^{\mathrm{cur}}_t=\nabla_S\log(\rho_t/g_N),\qquad
 I_H(t)=\int\norm{v^{\mathrm{cur}}_t}^2\rho_t\,\dd\omega_N,\qquad
 a_N(t)=\frac1N\E I_H(t)
\]
for the relative score, the Fisher information, and its normalized mean.
This is the notation of \cref{sec:entropy-clusters}; the superscript
distinguishes the score from the tree vectors $v_t$. Let $\mu_N(\dd t)=a_N(t)\,\dd t$. \Cref{lem:entropy-budget} implies that
$\mu_N([0,\infty))\le\beta^2/2$. By \eqref{eq:fisher-density-budget},
$\mu_{N_j}(I)\to\mu(I)$ for every
bounded interval $I$, where $\mu(\dd t)=a_F(t)\,\dd t$ is the finite measure
of \eqref{eq:entropy-measure}. We write $P^H_t$ for the diffusion
semigroup at fixed disorder and $\mathcal L_N$ for its generator.

\paragraph{Forks.}
For each fixed integer $K\ge1$, we use the product probability space
carrying $J$, the uniform initial point $X_0$, a full Brownian path
$(B_u)_{u\ge0}$, and $K$ full Brownian paths $B^{(1)},\dots,B^{(K)}$
on $[0,\infty)$. These coordinate blocks are mutually independent. We call
\[
 \zeta_N=\big(X_0,(B^{(b)})_{1\le b\le K},B\big)
\]
the fork seed, with its uniform-initial and Wiener product law; the
disorder has its independent standard Gaussian law. The seed space and law
depend on $K$, which we suppress in the notation. Fix $T\ge0$ and $L\ge0$. The trajectory $X$
solves \eqref{eq:hard-sde} driven by $B$. For $b\le K$, let $Z_b$ be the
solution of \eqref{eq:hard-sde} on $[0,T+L]$ with $Z_b(0)=X_0$, driven by
the Brownian motion that equals $B_u$ for $u\le T$ and
$B_T+B^{(b)}_{u-T}$ for $u\ge T$. Then $Z_b(u)=X_u$ for $u\le T$, and after
time $T$ the paths $Z_1,\dots,Z_K$ are continuations of $X$ from $X_T$ by
independent noises. We write
\[
 Y_b=Z_b(T+L),\qquad m_{N,T}=\frac1K\sum_{b=1}^KY_b,
\]
and call $m_{N,T}$ the center of the fork. When needed, we indicate the
dependence on $K$ and $L$ in the notation. The Brownian motion that drives
$Z_b$ is a standard Brownian motion independent of $(J,X_0)$. Hence each
pair $(J,Z_b)$ has the same law as $(J,(X_u)_{u\le T+L})$. The branches
share their past before $T$ and are not independent. Note that $\norm{m_{N,T}}\le\sqrt N$, because
$\norm{Y_b}=\sqrt N$.

\paragraph{Response moments.}
Let $U_b(t,u)$ be the variational response matrix of the path $Z_b$, defined as in \eqref{eq:response-matrix}. Each driving process is Brownian for its own full filtration and is independent of $(J,X_0)$, even though the branches share a common past. Therefore the proof of \eqref{eq:finite-lag-flow-moments} applies to every branch. With the constants $K_S$ and $\Theta_S$ defined there, its coordinate columns and normalized trace have fourth moment at most $K_S$ whenever $0\le t-u\le S$. In particular,
\[
 \E\left|\frac{\tr U_b(t,u)}N\right|^2\le\Theta_S^2,
 \qquad \E\norm{U_b(t,u)}_{\op}^2\le N\Theta_S^2.
\]
We use these estimates on the whole interval $[0,T+L]$, with $S=T+L$, and on the final interval $[T,T+L]$, with $S=L$. The latter bounds are uniform in $T$.

\begin{lemma}[Entropy comparison for two branches]\label{sel:fork-moments}
Let $T,a,d\ge0$, and let $Y$ and $Z$ be the endpoints of two branches at
time $T$ of lengths $a$ and $d$, that is, of two continuations of $X$ from
$X_T$ driven by independent Brownian motions. Set
$Q_{N,T}(a,d)=\ip{Y}{Z}/N$. Then, for $k\in\{1,2\}$,
\begin{equation}\label{sel:fork-moment-comparison}
 \left|\E Q_{N,T}(a,d)^k-\E C_N(T+2a+d,T+a)^k\right|
 \le2k\,e^{\beta p^2A_p(a+d)}\sqrt{a\,\mu_N([T,T+a])}+2e^{-c_pN}.
\end{equation}
Moreover, if $T_n\to\infty$ is a sequence along which the translated kernels
converge to a stationary cluster with correlation profile $c$, then, for
every $L<\infty$,
\begin{equation}\label{sel:fork-overlap-concentration}
 \lim_{n\to\infty}\limsup_{j\to\infty}\sup_{0\le a,d\le L}
 \E\left|Q_{N_j,T_n}(a,d)-c(a+d)\right|^2=0.
\end{equation}
\end{lemma}

\begin{proof}
For $k\in\{1,2\}$, let $\Phi_k(x)=(x/\sqrt N)^{\otimes k}$, with values in
$(\R^N)^{\otimes k}$ and its Euclidean norm, so that
$\ip{\Phi_k(x)}{\Phi_k(y)}=(\ip xy/N)^k$. On $S_N$, we have
$\norm{\Phi_k}=1$ and $\norm{D_S\Phi_k}_\op\le k/\sqrt N$. Fix the disorder.
For $0\le z\le a$, we define
\[
 f_z=P^H_{a-z}\Phi_k,\qquad g_z=P^H_{d+z}\Phi_k,\qquad
 F_J(z)=\int_{S_N}\rho_{T+z}\ip{f_z}{g_z}\,\dd\omega_N,
\]
where the semigroup acts on each coordinate of $\Phi_k$.

We first identify the endpoint values. By the Markov property and the
conditional independence of the two branches given the past up to time $T$,
we have $\E F_J(0)=\E Q_{N,T}(a,d)^k$. By the semigroup property,
$F_J(a)$ is the conditional expectation of
$\ip{\Phi_k(X_{T+a})}{\Phi_k(X_{T+2a+d})}$ given $J$, so that
$\E F_J(a)=\E C_N(T+2a+d,T+a)^k$.

Next we compute $F_J'$. We have $\partial_zf_z=-\mathcal L_Nf_z$, $\partial_zg_z=\mathcal L_Ng_z$, and
$\partial_z\rho_{T+z}=\mathcal L_N^*\rho_{T+z}$. By the product rule for the
generator, applied to each pair of coordinates,
$\mathcal L_N\ip{f}{g}=\ip{\mathcal L_Nf}{g}+\ip{f}{\mathcal L_Ng}
+2\sum_i\ip{D_Sf[e_i]}{D_Sg[e_i]}$, where $(e_i)$ is an orthonormal frame of
the tangent space. Using these identities, we find
\[
\begin{aligned}
 F_J'(z)&=\int\rho_{T+z}\left\{\mathcal L_N\ip{f_z}{g_z}
 -\ip{\mathcal L_Nf_z}{g_z}+\ip{f_z}{\mathcal L_Ng_z}\right\}\dd\omega_N\\
 &=2\int\rho_{T+z}\left\{\ip{f_z}{\mathcal L_Ng_z}
 +\sum_i\ip{D_Sf_z[e_i]}{D_Sg_z[e_i]}\right\}\dd\omega_N.
\end{aligned}
\]
When we apply \eqref{eq:gibbs-current-integration-by-parts} to each
coordinate of $f_z$ and $g_z$, the gradient terms cancel, and we obtain the
exact current identity
\begin{equation}\label{sel:fork-current-identity}
 F_J'(z)=-2\int\rho_{T+z}
 \ip{P^H_{a-z}\Phi_k}{D_S\big(P^H_{d+z}\Phi_k\big)\big[v^{\mathrm{cur}}_{T+z}\big]}
 \,\dd\omega_N.
\end{equation}
Note that this identity does not use stationarity of $\rho_{T+z}$: the
score $v^{\mathrm{cur}}_{T+z}$ records how far $\rho_{T+z}$ is from the
Gibbs density.

We now bound the right side of \eqref{sel:fork-current-identity} on the
event $G_N=\{Q_N\le A_p\}$. By Jensen's inequality,
$\norm{P^H_t\Phi_k}\le1$. In the proof of \cref{lem:waiting-clock-current},
the Bochner formula, the nonnegative Ricci curvature of $S_N$, and the Hessian bound of \eqref{eq:landscape-derivative-bounds} give the gradient
bound $\norm{\nabla_SP^H_t\varphi}^2\le e^{2\beta p^2Q_Nt}P^H_t\norm{\nabla_S\varphi}^2$
for smooth scalar $\varphi$; this is the classical gradient bound under a
lower bound on the curvature. Applied to $\varphi=\ip{w}{\Phi_k}$ for a unit
vector $w$ in $(\R^N)^{\otimes k}$, and followed by the supremum over $w$, it
gives
\[
 \norm{D_SP^H_t\Phi_k}_\op\le\frac k{\sqrt N}e^{\beta p^2A_pt}
 \qquad\text{on }G_N,
\]
with no factor depending on the dimension of $(\R^N)^{\otimes k}$. Using the
Cauchy--Schwarz inequality, we find that, on $G_N$ and for $z\le a$,
\[
 |F_J'(z)|\le\frac{2k}{\sqrt N}e^{\beta p^2A_p(a+d)}
 \int\rho_{T+z}\norm{v^{\mathrm{cur}}_{T+z}}\,\dd\omega_N
 \le2k\,e^{\beta p^2A_p(a+d)}\left(\frac{I_H(T+z)}N\right)^{1/2}.
\]
Taking expectations and using Jensen's inequality, then integrating in $z$
and using the Cauchy--Schwarz inequality, we obtain
\[
\begin{aligned}
 \int_0^a\E\big[\1_{G_N}|F_J'(z)|\big]\,\dd z
 &\le2k\,e^{\beta p^2A_p(a+d)}\int_0^a\sqrt{a_N(T+z)}\,\dd z\\
 &\le2k\,e^{\beta p^2A_p(a+d)}\sqrt{a\,\mu_N([T,T+a])}.
\end{aligned}
\]
On $G_N^c$, we have $|F_J(0)|,|F_J(a)|\le1$, and $\Prob(G_N^c)\le e^{-c_pN}$.
Writing
$\E[F_J(0)-F_J(a)]=-\E[\1_{G_N}\int_0^aF_J'(z)\,\dd z]+\E[\1_{G_N^c}(F_J(0)-F_J(a))]$
and combining the last two bounds with the identification of the endpoint
values, we obtain \eqref{sel:fork-moment-comparison}.

For \eqref{sel:fork-overlap-concentration}, fix $L$ and let $a,d\le L$. In
\eqref{sel:fork-moment-comparison}, we may replace $a\,\mu_N([T,T+a])$ by
$L\,\mu_N([T,T+L])$ and the exponential factor by $e^{2\beta p^2A_pL}$. Then
\eqref{eq:fisher-density-budget} implies that
$\mu_{N_j}([T,T+L])\to\mu([T,T+L])$ for each fixed $T$. By
\cref{prop:hard-bridge} on the horizon $T+3L$ and $|C_N|\le1$, the moments
$\E C_N(T+2a+d,T+a)^k$ converge to $C(T+2a+d,T+a)^k$ uniformly in
$a,d\le L$. Combining these facts with \eqref{sel:fork-moment-comparison},
we find
\[
 \limsup_{j\to\infty}\sup_{0\le a,d\le L}
 \left|\E Q_{N_j,T}(a,d)^k-C(T+2a+d,T+a)^k\right|
 \le4e^{2\beta p^2A_pL}\sqrt{L\,\mu([T,T+L])}.
\]
Along $T=T_n$, the right side tends to zero because $\mu$ is a finite
measure, and $C(T_n+2a+d,T_n+a)\to c(a+d)$ uniformly in $a,d\le L$, by the
locally uniform convergence of the translated correlations on
$[0,3L]^2$. Finally, writing $Q=Q_{N,T}(a,d)$ and $c_0=c(a+d)\in[0,1]$, we
have
\[
 \E|Q-c_0|^2=(\E Q^2-c_0^2)-2c_0(\E Q-c_0),
\]
and the convergence of the first two moments, uniformly in $a,d\le L$,
proves \eqref{sel:fork-overlap-concentration}.
\end{proof}

We now state the main result of this section. For $K\ge1$ and $L\ge0$, we
define
\begin{align}
 Q_{K,L}&=\frac1K+\left(1-\frac1K\right)c(2L),
 \label{sel:fork-norm}\\
 k_{K,L}(\tau)&=\frac{c(\tau)+(K-1)c(2L-\tau)}K,\qquad 0\le\tau\le L,
 \label{sel:fork-mixed-overlap}\\
 h_{K,L}&=e_\infty-\beta p\int_0^Lr(\tau)
 \left\{c(\tau)^{p-1}-k_{K,L}(\tau)^{p-1}\right\}\dd\tau.
 \label{sel:fork-center-energy}
\end{align}
Note that $Q_{K,L}$ is a convex combination of $1$ and $c(2L)$, so that
$q\le c(2L)\le Q_{K,L}\le1$, and that $0\le k_{K,L}\le1$.

\begin{proposition}[Finite forks and their center]\label{sel:finite-forks}
Fix $\beta>0$ and let $(c,r,\ell)$ be a stationary cluster with plateau
$q>0$ and energy $e_\infty$, realized along $T_n\to\infty$. Fix $K\ge1$ and
$L\ge0$, and let $m_{N,T}$ be the center of the fork at time $T$ with $K$
branches of length $L$.
\begin{enumerate}[label=(\alph*)]
\item\label{item:fork-limits}
The squared norm and the energy of the center satisfy
\begin{align}
 &\lim_{n\to\infty}\limsup_{j\to\infty}
 \E\left|\frac{\norm{m_{N_j,T_n}}^2}{N_j}-Q_{K,L}\right|^2=0,
 \label{sel:fork-center-concentration}\\
 &\lim_{n\to\infty}\limsup_{j\to\infty}
 \left|\frac{\E H_{N_j}(m_{N_j,T_n})}{N_j}-h_{K,L}\right|=0.
 \label{sel:fork-center-energy-limit}
\end{align}
\item\label{item:fork-ceiling}
The limiting energy is below the threshold:
\begin{equation}\label{sel:fork-center-threshold}
 h_{K,L}\le\Eth\,Q_{K,L}^{p/2}.
\end{equation}
\item\label{item:fork-lipschitz}
For every $T\ge0$ and $\varepsilon>0$, there exist algorithms
$\mathcal A^\varepsilon_{N,T}=\mathcal A^\varepsilon_{N,T}(J,\zeta_N)$ with
values in $B_N$ and seed $\zeta_N$, which are $L_{\mathrm{Lip}}$-Lipschitz for a
constant $L_{\mathrm{Lip}}$ depending only on $p$, $\beta$, $T$, $K$, $L$, and
$\varepsilon$, such that
\begin{equation}\label{sel:fork-lipschitz}
 \limsup_{j\to\infty}\frac1{N_j}\E\norm{\mathcal A^\varepsilon_{N_j,T}
 -\Pi_{B_{N_j}}\big(m_{N_j,T}/\sqrt{Q_{K,L}}\big)}^2\le\varepsilon^2.
\end{equation}
Moreover, along $T=T_n$, the vector $\Pi_{B_N}(m_{N,T}/\sqrt{Q_{K,L}})$ may
be replaced by $m_{N,T}/\sqrt{Q_{K,L}}$, or by
$\widehat m_{N,T}=\sqrt N\,m_{N,T}/\norm{m_{N,T}}$ (with $\widehat m_{N,T}=0$
if $m_{N,T}=0$): the scaled mean-square distance between
$\Pi_{B_N}(m_{N,T}/\sqrt{Q_{K,L}})$ and either vector tends to zero as
$j\to\infty$ and then $n\to\infty$.
\end{enumerate}
\end{proposition}

By \eqref{sel:tensor-ui}, the normalized energies of the algorithms in
part~\ref{item:fork-lipschitz} are bounded by $Q_N$. Part~\ref{item:fork-lipschitz}
also transfers overlaps. If $V_N$ is a random vector on the same probability
space with $\norm{V_N}\le\sqrt N$ (for example $V_N=X_s$ with $s\le T$), then, by the Cauchy--Schwarz inequality, we have
\begin{equation}\label{eq:fork-overlap-transfer}
 \left|\E\frac{\ip{\mathcal A^\varepsilon_{N,T}}{V_N}}N
 -\E\frac{\ip{m_{N,T}}{V_N}}{N\sqrt{Q_{K,L}}}\right|
 \le\left(\frac1N\E\norm{\mathcal A^\varepsilon_{N,T}
 -m_{N,T}/\sqrt{Q_{K,L}}}^2\right)^{1/2}.
\end{equation}

\begin{proof}
The proof has five steps. In the first two, we compute the energy of the
center exactly by Gaussian integration by parts and subtract the energy of
a single branch. In the third, we show that the common past before $T$
contributes nothing in the dimension limit, and in the fourth we evaluate
the contribution of the last $L$ units of time; this proves
part~\ref{item:fork-limits}. In the fifth step, we construct the Lipschitz
algorithms of part~\ref{item:fork-lipschitz} and deduce
part~\ref{item:fork-ceiling} from the threshold. We write $t=T+L$.

\paragraph{Step 1: Gaussian integration by parts.}
For an ordered index $I=(i_1,\dots,i_p)$, let
$G_I(x)=N^{-(p-1)/2}x_{i_1}\cdots x_{i_p}$, so that
$H_N=\sum_IJ_IG_I$. Since $\sum_IG_I(x)G_I(z)=N(\ip xz/N)^p$,
differentiating first in $x_i$ and then in $z_l$ gives
\begin{equation}\label{sel:fork-covariance-derivative}
 \sum_I\partial_iG_I(x)\,\partial_lG_I(z)
 =p\left(\frac{\ip xz}N\right)^{p-1}\delta_{il}
 +\frac{p(p-1)}N\left(\frac{\ip xz}N\right)^{p-2}z_ix_l.
\end{equation}
Differentiating the actual terminal solution into $L^2$ with respect to the disorder, as in the proof of
\cref{prop:energy-bridge}, we have
\[
 \partial_{J_I}m_{N,T}
 =\frac\beta K\sum_{b=1}^K\int_0^tU_b(t,u)\nabla G_I(Z_b(u))\,\dd u.
\]
The seed is independent of $J$, and the response has the column and operator moments above. As in the proof of \cref{prop:energy-bridge}, apply Gaussian integration by parts to the expected coordinate observable in each disorder coefficient, using its genuine $L^2$ derivative:
\[
 \E\big[J_IG_I(m_{N,T})\big]=\E\big[\ip{\nabla G_I(m_{N,T})}{\partial_{J_I}m_{N,T}}\big],
\]
and $\E H_N(m_{N,T})$ is the sum of these terms over $I$. For $u\le t$, we
write
\[
 \gamma_b(u)=\frac{\ip{Y_b}{Z_b(u)}}N,\qquad
 \bar\gamma_b(u)=\frac{\ip{m_{N,T}}{Z_b(u)}}N
 =\frac1K\sum_{b'=1}^K\frac{\ip{Y_{b'}}{Z_b(u)}}N.
\]
Contracting \eqref{sel:fork-covariance-derivative} with $U_b(t,u)$ gives the
exact identity
\begin{align}
 \frac{\E H_N(m_{N,T})}N
 ={}&\frac{\beta p}K\sum_{b=1}^K\int_0^t
 \E\left[\frac{\tr U_b(t,u)}N\,\bar\gamma_b(u)^{p-1}\right]\dd u
 +\mathcal R_{N,T},\label{sel:fork-exact-ibp}\\
 \mathcal R_{N,T}={}&\frac{\beta p(p-1)}K\sum_{b=1}^K\int_0^t
 \E\left[\bar\gamma_b(u)^{p-2}
 \frac{\ip{Z_b(u)}{U_b(t,u)m_{N,T}}}{N^2}\right]\dd u.\notag
\end{align}
Using $\norm{m_{N,T}},\norm{Z_b(u)}\le\sqrt N$ and $|\bar\gamma_b|\le1$, we
find that $|\mathcal R_{N,T}|\le\beta p(p-1)t\,\Theta_t/\sqrt N$.

\paragraph{Step 2: the exact subtraction.}
The same computation with $m_{N,T}$ replaced by $Y_b$ gives
$\E H_N(Y_b)/N=\beta p\int_0^t\E[N^{-1}\tr U_b(t,u)\gamma_b(u)^{p-1}]\,\dd u
+O(t\,\Theta_t/\sqrt N)$. The pair $(J,Y_b)$ has the law of $(J,X_t)$, and
therefore $\E H_N(Y_b)/N=e_N(t)$, where $e_N(t)=\E H_N(X_t)/N$. Averaging over $b$ and
subtracting \eqref{sel:fork-exact-ibp}, before taking any limit, we obtain
\begin{equation}\label{sel:fork-exact-subtraction}
\begin{aligned}
 e_N(t)-\frac{\E H_N(m_{N,T})}N
 ={}&\frac{\beta p}K\sum_{b=1}^K\int_0^t
 \E\left[\frac{\tr U_b(t,u)}N
 \big\{\gamma_b(u)^{p-1}-\bar\gamma_b(u)^{p-1}\big\}\right]\dd u\\
 &+O\big(t\,\Theta_tN^{-1/2}\big),
\end{aligned}
\end{equation}
where the implicit constant depends only on $p$ and $\beta$.
Note that this identity retains the whole history of the trajectory.

\paragraph{Step 3: the common past.}
Let $u\le T$. Then $Z_b(u)=X_u$ for every $b$, and the pair $(Y_b,X_u)$ has
the law of $(X_t,X_u)$. By \cref{prop:hard-bridge} and $|C_N|\le1$, we
have $\E|\gamma_b(u)-C(t,u)|^2\to0$ as $N\to\infty$, uniformly in $u\le T$.
The same holds for $\bar\gamma_b(u)$, which is the average of
$\gamma_{b'}(u)$ over $b'$. The function $x\mapsto x^{p-1}$ is
$(p-1)$-Lipschitz on $[-1,1]$, and the normalized trace has second moment at most $\Theta_t^2$. Using the
Cauchy--Schwarz inequality, we find that the part of the integral in
\eqref{sel:fork-exact-subtraction} over $u\le T$ is at most
\[
 \beta p(p-1)\,T\,\Theta_t
 \sup_{b\le K,\,u\le T}\big(\E|\gamma_b(u)-\bar\gamma_b(u)|^2\big)^{1/2}
\]
in absolute value. This tends to zero as $N\to\infty$ with $T$ fixed. We
use no estimate on the decay of the memory here. The contribution of the
past is small because, at a fixed time horizon, both overlaps concentrate
around the same deterministic value $C(t,u)$.

\paragraph{Step 4: the last $L$ units of time.}
Let $u=T+v$ with $0\le v\le L$, and write $\tau=L-v$. The pair
$(Y_b,Z_b(T+v))$ has the law of $(X_t,X_{T+v})$. By \cref{prop:hard-bridge},
the overlap $\gamma_b(T+v)$ converges in $L^2$ to $C(T+L,T+v)$ as
$N\to\infty$, uniformly in $v$. Along $T=T_n$, this tends to $c(L-v)=c(\tau)$, uniformly in
$v\le L$. For $b'\ne b$, the vectors $Y_{b'}$ and $Z_b(T+v)$ are the
endpoints of two branches at $T$ of lengths $L$ and $v$. By
\eqref{sel:fork-overlap-concentration}, the overlap
$\ip{Y_{b'}}{Z_b(T+v)}/N$ converges in $L^2$ to $c(L+v)=c(2L-\tau)$, in the
double limit $j\to\infty$ and then $n\to\infty$, uniformly in $v\le L$.
Averaging over $b'$, we find
\[
 \lim_{n\to\infty}\limsup_{j\to\infty}\sup_{0\le v\le L}
 \Big(\E|\gamma_b(T_n+v)-c(L-v)|^2
 +\E|\bar\gamma_b(T_n+v)-k_{K,L}(L-v)|^2\Big)=0.
\]
On this interval, we have $t-u\le L$, and the normalized response trace has second moment at most $\Theta_L^2$, uniformly in $T$. As in Step~3, using the Cauchy--Schwarz inequality, we may replace $\gamma_b^{p-1}-\bar\gamma_b^{p-1}$ in
\eqref{sel:fork-exact-subtraction} by the deterministic weight
\[
 \omega(\tau)=c(\tau)^{p-1}-k_{K,L}(\tau)^{p-1},\qquad\tau=t-u,
\]
with an error whose $\limsup_j$ tends to zero as $n\to\infty$.

It remains to pass to the limit in
$\beta p\int_T^t\bar r_N(t,u)\,\omega(t-u)\,\dd u$, where
$\bar r_N(t,u)=\E\tr U_b(t,u)/N$ is the averaged response of
\cref{lem:brown-force-response}, which does not depend on $b$.
\Cref{lem:brown-force-response} and \cref{prop:hard-bridge} imply that
$\int_{s_1}^{s_2}\bar r_N(t,u)\,\dd u\to\int_{s_1}^{s_2}R(t,u)\,\dd u$
for all $T\le s_1\le s_2\le t$. The signed measure $\bar r_N(t,u)\,\dd u$ on
$[T,t]$ has total variation
\[
 \int_T^t|\bar r_N(t,u)|\,\dd u
 \le\int_T^t\E\left|\frac{\tr U_b(t,u)}N\right|\,\dd u
 \le L\,\Theta_L.
\]
As $R\ge0$, the integral of $R(t,\cdot)$ over $[T,t]$ is the limit of the
masses of these signed measures and obeys the same bound. The weight $\omega$ is continuous on $[0,L]$ and
is the uniform limit of step functions. For a step function, the integral converges as $N\to\infty$, and the total variation bound controls the error
of the approximation, uniformly in $N$. We conclude that, for fixed $T$,
\[
 \lim_{N\to\infty}\int_T^t\bar r_N(t,u)\,\omega(t-u)\,\dd u
 =\int_T^tR(t,u)\,\omega(t-u)\,\dd u.
\]
Along $T=T_n$, we have $R(T_n+L,T_n+v)\to r(L-v)$ uniformly in $v\le L$, by
the translated convergence on the closed causal triangle. Moreover,
$e_N(t)\to e_\beta(T+L)$ by \cref{prop:energy-bridge}, and
$e_\beta(T_n+L)\to e_\infty$. Inserting these limits into
\eqref{sel:fork-exact-subtraction}, we obtain
\[
 \lim_{n\to\infty}\limsup_{j\to\infty}
 \left|\frac{\E H_{N_j}(m_{N_j,T_n})}{N_j}
 -e_\infty+\beta p\int_0^Lr(\tau)\,\omega(\tau)\,\dd\tau\right|=0,
\]
which is \eqref{sel:fork-center-energy-limit} by the definition
\eqref{sel:fork-center-energy}.

For the norm, we have
\[
 \frac{\norm{m_{N,T}}^2}N=\frac1K+\frac1{K^2}\sum_{b\ne b'}
 \frac{\ip{Y_b}{Y_{b'}}}N.
\]
Each overlap $\ip{Y_b}{Y_{b'}}/N$ with $b\ne b'$ has the law of
$Q_{N,T}(L,L)$, which tends to $c(2L)$ in $L^2$ in the double limit, by
\eqref{sel:fork-overlap-concentration}. There are $K(K-1)$ such terms, and
$K(K-1)/K^2=1-1/K$. This proves \eqref{sel:fork-center-concentration} and
completes the proof of part~\ref{item:fork-limits}.

\paragraph{Step 5: Lipschitz algorithms and the threshold.}
Fix $T$ and $\varepsilon>0$. We may assume that $\varepsilon<1$. Let
$\varepsilon_1=\varepsilon/(2\sqrt K)$. We use the hard--soft approximation
of \cite[Lemmas~3.1--3.2]{sellke2024}. We first check its hypotheses.
In \cite{sellke2024}, the Hamiltonian is our $H_N$, with an ordered tensor of
independent standard Gaussian coefficients and the Euclidean distance between
coefficient tensors, and the dynamics is
$\dd x_v=(\beta_{\mathrm S}\nabla_{\mathrm{sp}}H_N(x_v)
-(N-1)(2N)^{-1}x_v)\,\dd v+P_{x_v}\,\dd W_v$, with noise coefficient one.
With $\beta_{\mathrm S}=\beta/2$, $X_u=x_{2u}$, and $B_u=W_{2u}/\sqrt2$, this
is \eqref{eq:hard-sde}, as in the proof of \cref{prop:hard-bridge}, and our
horizon $t$ becomes $2t$. Lemma~3.1 there states that, for a fixed initial
point on $S_N$ and a Hamiltonian satisfying their derivative bounds (which
hold for $H_N$ with probability at least $1-e^{-c_1N}$, by
\cite[Proposition~1.11]{sellke2024}), the hard dynamics stays within
$\varepsilon_1\sqrt N$ of a soft spherical dynamics driven by the same
Brownian motion, on the whole horizon, with probability at least
$1-e^{-c_1N}$. Here the confinement parameter of the soft dynamics is chosen large, depending on $p$, $\beta$, $t$, and $\varepsilon_1$, and $c_1>0$ and the
threshold for $N$ depend on the same parameters. Lemma~3.2 there states that, for a fixed initial point and
almost every Brownian path, there exists a map $\ell$ from coefficient tensors
to $B_N$, Lipschitz with a constant depending only on $p$, $\beta$, $t$, and
$\varepsilon_1$, such that $(1+\varepsilon_1)\ell(J)$ equals the soft
endpoint with probability at least $1-e^{-c_1N}$, where the probability is
over both the disorder and the Brownian motion. The coefficient-Lipschitz extension uses Kirszbraun's theorem in the proof of
Lemma~3.2 of \cite{sellke2024}; footnote~2 there gives a constructive choice
that depends measurably on the Brownian path.

The construction of \cite{sellke2024} is for a fixed initial point. To
include the uniform initial point, we choose measurably an orthogonal matrix
$O_{x_0}$ such that $O_{x_0}(\sqrt N e_1)=x_0$. Rotating the tensor by
$(O_{x_0}^{\mathsf T})^{\otimes p}$ and the Brownian path by
$O_{x_0}^{\mathsf T}$, solving from $\sqrt Ne_1$, and rotating the output back
by $O_{x_0}$ gives the solution from $x_0$, because $H_N$ and
\eqref{eq:hard-sde} are equivariant under rotations. Conditionally on
$X_0=x_0$, the rotated tensor and Brownian path have the same joint law as
the original ones, and the tensor rotation is an isometry of
$(\R^N)^{\otimes p}$, which does not change Lipschitz constants.

We apply this to each concatenated path $Z_b$, whose driving Brownian
motion is a function of the seed $\zeta_N$. We obtain maps
$\ell_b(J)=\ell_b(J,\zeta_N)\in B_N$, $b\le K$, which are Lipschitz in $J$
with a constant $L_{\mathrm{Lip},1}$ depending only on $p$, $\beta$, $t$, and
$\varepsilon_1$. Let $\mathcal G_N$ be the event on which $H_N$ satisfies the
derivative bounds of \cite{sellke2024} and, for every $b\le K$, the
conclusions of both lemmas hold for the branch $b$. On $\mathcal G_N$, we
have $\norm{(1+\varepsilon_1)\ell_b(J)-Y_b}\le\varepsilon_1\sqrt N$ for every
$b\le K$, and a union bound gives $\Prob(\mathcal G_N^c)\le3Ke^{-c_1N}$. The
union bound does not require the branches to be independent. We define
\begin{equation}\label{sel:fork-projected-algorithm}
 \mathcal A^\varepsilon_{N,T}(J)
 =\Pi_{B_N}\left(\frac{1+\varepsilon_1}{\sqrt{Q_{K,L}}}
 \frac1K\sum_{b=1}^K\ell_b(J)\right).
\end{equation}
Since $\Pi_{B_N}$ is $1$-Lipschitz and $Q_{K,L}\ge1/K$, this is an
$L_{\mathrm{Lip}}$-Lipschitz algorithm with seed $\zeta_N$ and
$L_{\mathrm{Lip}}=2L_{\mathrm{Lip},1}\sqrt K$. On $\mathcal G_N$, the distance between
$\mathcal A^\varepsilon_{N,T}$ and $\Pi_{B_N}(m_{N,T}/\sqrt{Q_{K,L}})$ is at
most $\varepsilon_1\sqrt N/\sqrt{Q_{K,L}}\le\varepsilon_1\sqrt{KN}
=\varepsilon\sqrt N/2$; on
$\mathcal G_N^c$, it is at most $2\sqrt N$. Using
$\Prob(\mathcal G_N^c)\le3Ke^{-c_1N}$, we find
\[
 \frac1N\E\norm{\mathcal A^\varepsilon_{N,T}
 -\Pi_{B_N}(m_{N,T}/\sqrt{Q_{K,L}})}^2
 \le\frac{\varepsilon^2}4+12Ke^{-c_1N},
\]
which proves \eqref{sel:fork-lipschitz}.

Next we compare the three normalizations. Let
$\nu_N=\norm{m_{N,T}}^2/N$. For $x\ge0$, we have
\begin{equation}\label{eq:sqrt-comparison}
 \left|\sqrt{x/Q_{K,L}}-1\right|
 =\frac{|x-Q_{K,L}|}{\sqrt{Q_{K,L}}\,(\sqrt x+\sqrt{Q_{K,L}})}
 \le\frac{|x-Q_{K,L}|}{Q_{K,L}}.
\end{equation}
The scaled distance between $m_{N,T}/\sqrt{Q_{K,L}}$ and its projection
onto $B_N$ is $(\sqrt{\nu_N/Q_{K,L}}-1)_+$. If $m_{N,T}\ne0$, the scaled
distance between the projection and $\widehat m_{N,T}$ is
$(1-\sqrt{\nu_N/Q_{K,L}})_+$, and both vectors vanish if $m_{N,T}=0$. By
\eqref{eq:sqrt-comparison} and $Q_{K,L}\ge q$, both distances are at most
$|\nu_N-Q_{K,L}|/q$. Then \eqref{sel:fork-center-concentration} implies that
their mean squares tend to zero as $j\to\infty$ and then $n\to\infty$. This completes
the proof of part~\ref{item:fork-lipschitz}.

Finally, we prove part~\ref{item:fork-ceiling}. Since
$\mathcal A^\varepsilon_{N,T}$ is an $L_{\mathrm{Lip}}$-Lipschitz algorithm with seed
$\zeta_N$, \cref{sel:absolute-field-cost}, applied with any
$\zeta_N$-measurable direction and with the nonnegative overlap term
dropped, gives
\[
 \limsup_{N\to\infty}\frac{\E H_N(\mathcal A^\varepsilon_{N,T})}N\le\Eth.
\]
This is the pure-model threshold, the case $\kappa=0$ of
\cref{sel:published-field-threshold}, averaged over the seed. Next we
compare energies. Note that $m_{N,T}/\sqrt{Q_{K,L}}$ has norm at most
$\sqrt{N/q}$. Using \eqref{eq:general-energy-transfer} with $R=q^{-1/2}$ and
the homogeneity $H_N(m_{N,T}/\sqrt{Q_{K,L}})=Q_{K,L}^{-p/2}H_N(m_{N,T})$, we
find
\[
 Q_{K,L}^{-p/2}\frac{\E H_N(m_{N,T})}N
 \le\frac{\E H_N(\mathcal A^\varepsilon_{N,T})}N
 +p\,q^{-(p-1)/2}Q_*\,\Delta_{N,T},
\]
where
\[
 \Delta_{N,T}=\left(\frac1N\E\norm{\mathcal A^\varepsilon_{N,T}
 -m_{N,T}/\sqrt{Q_{K,L}}}^2\right)^{1/2}.
\]
By \eqref{sel:fork-lipschitz} and the comparison of normalizations, we have
\[
 \limsup_{n\to\infty}\limsup_{j\to\infty}\Delta_{N_j,T_n}\le\varepsilon.
\]
We take
$j\to\infty$ with $T=T_n$ fixed, then $n\to\infty$, and finally
$\varepsilon\downarrow0$, in the order \eqref{sel:aging-limit-order}. By
\eqref{sel:fork-center-energy-limit}, we obtain
$Q_{K,L}^{-p/2}h_{K,L}\le\Eth$, which is \eqref{sel:fork-center-threshold}.
The constants of the algorithms depend on $T$, but we use them only at a
fixed $T$, before the limit $n\to\infty$.
\end{proof}

We next record the limits of the fork parameters as $K\to\infty$ and then
$L\to\infty$. Recall that $d_L=c(2L)$ and $h_L$ are given by
\eqref{sel:prefix-norm} and \eqref{sel:prefix-energy}, with the cluster
energy $e_\infty$, and that $h$ is the quantity in
\eqref{sel:limiting-center-energy}.

\begin{lemma}[The fork parameters as $K\to\infty$ and $L\to\infty$]
\label{lem:fork-parameter-limits}
Let $(c,r,\ell)$ be a stationary cluster with plateau $q\ge0$ and energy
$e_\infty$.
\begin{enumerate}[label=(\roman*)]
\item\label{item:fork-K-limit}
For every $L\ge0$ and $K\ge1$,
\begin{equation}\label{eq:fork-K-limit}
 0\le Q_{K,L}-d_L\le\frac1K,\qquad
 |h_{K,L}-h_L|\le\frac{\beta p(p-1)}K.
\end{equation}
\item\label{item:fork-L-limit}
As $L\to\infty$, we have $d_L\to q$ and
$h_L\to h=e_\infty-\beta\{1-q^p-pq^{p-1}(1-q)\}$.
\item\label{item:fork-saturation}
If $q>0$ and $\beta^2p(p-1)q^{p-2}(1-q)^2=1$, then
\begin{equation}\label{sel:fork-saturation}
 \lim_{L\to\infty}\lim_{K\to\infty}\frac{h_{K,L}}{Q_{K,L}^{p/2}}=\Eth.
\end{equation}
\end{enumerate}
\end{lemma}

\begin{proof}
For part~\ref{item:fork-K-limit}, we have
$Q_{K,L}-d_L=(1-c(2L))/K\in[0,1/K]$. By \eqref{sel:prefix-energy} and
\eqref{sel:fork-center-energy},
\[
 h_{K,L}-h_L=\beta p\int_0^Lr(\tau)
 \left\{k_{K,L}(\tau)^{p-1}-c(2L-\tau)^{p-1}\right\}\dd\tau.
\]
The integrand is dominated by $r(\tau)$ times a quantity that tends to zero
uniformly. Indeed, $k_{K,L}(\tau)-c(2L-\tau)=(c(\tau)-c(2L-\tau))/K$, and both
$k_{K,L}(\tau)$ and $c(2L-\tau)$ lie in $[0,1]$. By the mean value theorem,
\[
 \left|k_{K,L}(\tau)^{p-1}-c(2L-\tau)^{p-1}\right|\le\frac{p-1}K.
\]
Together with $\int_0^Lr\le1$, this proves \eqref{eq:fork-K-limit}.

Part~\ref{item:fork-L-limit} is part of \cref{lem:center-energy}. For
completeness, we recall the computation. Using $r=-c'$, we find that
$\int_0^Lr(\tau)c(\tau)^{p-1}\,\dd\tau=(1-c(L)^p)/p$. As $c$ is
nonincreasing, we have $c(2L-\tau)\in[q,c(L)]$ for $0\le\tau\le L$, and $\int_0^Lr=1-c(L)$. Thus the integral of
$r(\tau)c(2L-\tau)^{p-1}$ over $[0,L]$ lies between $q^{p-1}(1-c(L))$ and $c(L)^{p-1}(1-c(L))$. Both bounds
tend to $q^{p-1}(1-q)$. Inserting these two integrals into
\eqref{sel:prefix-energy}, we find that $h_L$ tends to
$e_\infty+\beta pq^{p-1}(1-q)-\beta(1-q^p)=h$.

For part~\ref{item:fork-saturation}, part~\ref{item:fork-K-limit} and
$d_L\ge q>0$ imply that $h_{K,L}/Q_{K,L}^{p/2}\to h_L/d_L^{p/2}$ as $K\to\infty$,
and part~\ref{item:fork-L-limit} gives $h_L/d_L^{p/2}\to h/q^{p/2}$ as
$L\to\infty$. By the equality case of \eqref{sel:center-amgm} in
\cref{lem:center-energy}, the marginality hypothesis implies that
$h/q^{p/2}=\Eth$.
\end{proof}

Part~\ref{item:fork-saturation} uses marginality, which is proved in
\cref{sel:positive-marginality} from part~\ref{item:threshold-ceiling} of
\cref{prop:threshold-inputs}. The proof of that part in
\cref{sec:general-proof-inputs} uses only \cref{sel:finite-forks} and
part~\ref{item:fork-K-limit}, and there is no circularity.

\section{Adaptive disclosure and orthogonality to early queries}
\label{sec:general-disclosure}

The overlap penalty of \cref{sel:absolute-field-cost} applies to directions
independent of $J$. At positive times, however, an earlier state $X_s$ depends on $J$. The goal of this section is to prove that an algorithm that saturates the threshold is almost
orthogonal to every subspace explored by a fixed finite number of gradient
queries. We condition $J$ on the information revealed by the queries. This is adaptive Gaussian conditioning, in the line of Bolthausen's
iterative construction of solutions of the TAP equations
\cite[Section~3]{bolthausen2014}. Bayati and Montanari used the same method
in their analysis of approximate message passing
\cite[arXiv version~4, Section~3.8, Lemma~11]{bayatimontanari2011}, and
they credit it to Bolthausen. The underlying facts on Gaussian vectors and
their orthogonal decompositions are classical; see \cite[Chapter~9, Theorem~9.1, pp.~127--128, and Theorems~9.8--9.9, pp.~131--132]{janson1997}. We use projection onto the span of the deterministic observations, which also permits linearly dependent observation families.
El Alaoui, Montanari, and Sellke conditioned the disorder tensor on adaptive
gradient observations in their analysis of message passing for spin
glasses \cite[arXiv version~1, Appendix~A, Lemma~A.2, and Theorem~7(a)]{eams2021}, and Huang and Sellke extended this conditioning to
multi-species spherical models \cite[arXiv version~1, Appendix~A, Lemma~A.2, and Theorem~3(a)]{hs2024amp}. The displayed inverse-Gram formulas in these AMP analyses use rank and conditioning hypotheses. Our projection-based disclosure below also covers dependent queries. Here we disclose all tensor blocks needed to
determine gradients throughout the explored span. We combine this disclosure with the external-field threshold to prove forgetting of a fixed
finite history.

\paragraph{Outer parameters.}
Let $\theta$ be a finite list of outer parameters with a prescribed order of
limits; in \cref{sec:general-transfer}, $\theta=(n,\varepsilon,K,L)$, in the
order \eqref{sel:aging-limit-order}. For nonnegative quantities
$F_{N,\theta}$, we write $\lim_\theta\limsup_NF_{N,\theta}=0$ if the upper
limit in $N$, followed by the upper limits in the outer parameters in their
prescribed order, vanishes. We do not assume that any intermediate limit
exists.

\paragraph{Standing hypotheses.}
For each $N$ and $\theta$, let $\zeta_{N,\theta}$ be a random variable independent of $J$, which contains all
auxiliary randomness used both by the algorithms and by the disclosure
procedure below. Its measurable space and law may depend on $\theta$;
below we write $\zeta_N$ with $\theta$ suppressed. For each $N$ and $\theta$, let
$\mathcal A_{N,\theta}=\mathcal A_{N,\theta}(J,\zeta_N)$ be an
$L_{\mathrm{Lip},\theta}$-Lipschitz algorithm with seed $\zeta_N$. The constant
$L_{\mathrm{Lip},\theta}$ is deterministic and independent of $N$, but it may be
unbounded in the outer limits. We define the energy deficit
\begin{equation}\label{sel:adaptive-deficit}
 \delta_\theta=\max\left\{0,\;\Eth-\liminf_{N\to\infty}
 \E\frac{H_N(\mathcal A_{N,\theta})}N\right\}.
\end{equation}
\Cref{sel:absolute-field-cost} implies that
$\limsup_N\E H_N(\mathcal A_{N,\theta})/N\le\Eth$. The hypothesis
$\lim_\theta\delta_\theta=0$ says exactly that the algorithms saturate the
threshold in mean.

For a subspace $W\subset\R^N$, we write
\[
 \mathcal T(W)=(W^\perp)^{\otimes p}\subset(\R^N)^{\otimes p},
\]
and $P_V$ for the orthogonal projection onto a subspace $V$ of $\R^N$ or of
$(\R^N)^{\otimes p}$. If $W\subset W'$, then $\mathcal T(W')\subset\mathcal T(W)$,
and we write $\mathcal T(W)\ominus\mathcal T(W')$ for the orthogonal
complement of $\mathcal T(W')$ in $\mathcal T(W)$. We choose orthonormal
bases of measurable subspaces measurably, by Gram--Schmidt in a fixed order
with deterministic fallbacks.

\begin{lemma}[Adaptive tensor disclosure]\label{sel:adaptive-disclosure}
Let $W_0\subset\R^N$ be a $\zeta_N$-measurable subspace, and let
$\mathcal F_0$ be the $\sigma$-algebra generated by $\zeta_N$ and
$P_{\mathcal T(W_0)^\perp}J$. For $j\ge0$, given $W_j$ and $\mathcal F_j$,
let $u_j$ be an $\mathcal F_j$-measurable vector that is either zero or a
unit vector in $W_j^\perp$, and set
\[
 W_{j+1}=W_j+\Span(u_j),\qquad
 \mathcal F_{j+1}=\sigma\Big(\mathcal F_j,\;
 P_{\mathcal T(W_j)\ominus\mathcal T(W_{j+1})}J\Big).
\]
Then, for every $j$:
\begin{enumerate}[label=(\roman*)]
\item\label{item:disclosure-gaussian}
$P_{\mathcal T(W_j)^\perp}J$ is $\mathcal F_j$-measurable, and conditionally
on $\mathcal F_j$, the tensor $P_{\mathcal T(W_j)}J$ is standard Gaussian on
$\mathcal T(W_j)$;
\item\label{item:disclosure-gradient}
if $x_j\in W_j$ is $\mathcal F_j$-measurable, then $\nabla H_N(x_j)$ is
$\mathcal F_j$-measurable. In particular, $u_j$ may be chosen as the
normalized projection of $\nabla H_N(x_j)$ onto $W_j^\perp$, and as zero if
that projection vanishes.
\end{enumerate}
\end{lemma}

\begin{proof}
The conditional Gaussian assertions are understood almost surely with respect to the disclosed history. A zero query discloses the zero tensor block and leaves the residual law unchanged.
We prove part~\ref{item:disclosure-gaussian} by induction on $j$. For
$j=0$, conditionally on $\zeta_N$, the tensor $J$ is standard Gaussian,
because it is independent of $\zeta_N$, and the subspaces $\mathcal T(W_0)$
and $\mathcal T(W_0)^\perp$ are fixed. The projections of a standard
Gaussian vector onto two complementary orthogonal subspaces are independent
standard Gaussian vectors on them. Revealing $P_{\mathcal T(W_0)^\perp}J$
leaves $P_{\mathcal T(W_0)}J$ standard Gaussian on
$\mathcal T(W_0)$.

Suppose that part~\ref{item:disclosure-gaussian} holds for $j$. Conditionally on
$\mathcal F_j$, the vector $u_j$ and the subspaces in the orthogonal
decomposition
\[
 \mathcal T(W_j)=\mathcal T(W_{j+1})\oplus
 \big(\mathcal T(W_j)\ominus\mathcal T(W_{j+1})\big)
\]
are fixed, and $P_{\mathcal T(W_j)}J$ is standard Gaussian on
$\mathcal T(W_j)$. Its projections onto the two summands are
independent standard Gaussian vectors. Revealing the second leaves the
first, which is $P_{\mathcal T(W_{j+1})}J$, standard Gaussian on
$\mathcal T(W_{j+1})$. Since
$\mathcal T(W_{j+1})^\perp=\mathcal T(W_j)^\perp\oplus
(\mathcal T(W_j)\ominus\mathcal T(W_{j+1}))$, the projection
$P_{\mathcal T(W_{j+1})^\perp}J$ is $\mathcal F_{j+1}$-measurable. This
completes the induction. At each step, we choose the new direction before we reveal the
corresponding part of $J$; this is what makes the induction work, and we
never condition on a subspace that depends on undisclosed randomness.

For part~\ref{item:disclosure-gradient}, let $x\in W_j$ and $y\in\R^N$. We
have
\[
 \ip{\nabla H_N(x)}{y}=N^{-(p-1)/2}\Big\langle J,\;\sum_{i=1}^p
 x^{\otimes(i-1)}\otimes y\otimes x^{\otimes(p-i)}\Big\rangle.
\]
Since $p\ge2$, each tensor in the sum has at least one factor $x\in W_j$.
It is orthogonal to every product tensor with all factors in $W_j^\perp$,
and these product tensors span $\mathcal T(W_j)$. The right side is equal to
$N^{-(p-1)/2}$ times the inner product of $P_{\mathcal T(W_j)^\perp}J$ with
that tensor. By part~\ref{item:disclosure-gaussian}, it is
$\mathcal F_j$-measurable when $x=x_j$ is.
\end{proof}

\begin{proposition}[Orthogonality to every fixed finite query span]
\label{sel:adaptive-orthogonality}
Assume the standing hypotheses and $\lim_\theta\delta_\theta=0$. Fix integers
$k_0\ge0$ and $M\ge0$. For each $N$ and $\theta$, let
$W_0\subset W_1\subset\dots\subset W_M$ be subspaces constructed as in
\cref{sel:adaptive-disclosure} from the seed $\zeta_N$, with
$\dim W_0\le k_0$. Then
\begin{equation}\label{sel:finite-span-orthogonality}
 \lim_\theta\limsup_{N\to\infty}
 \E\frac{\norm{P_{W_M}\mathcal A_{N,\theta}}^2}N=0.
\end{equation}
In particular, if $Y_N\in W_M\cap B_N$ is a random vector, then
$\lim_\theta\limsup_N\E|\ip{\mathcal A_{N,\theta}}{Y_N}/N|^2=0$.
\end{proposition}

\begin{proof}
For $0\le j\le M$, we define
\[
 \psi_{j,\theta}=\limsup_{N\to\infty}
 \E\frac{\norm{P_{W_j}\mathcal A_{N,\theta}}^2}N.
\]
We bound $\psi_{0,\theta}$ directly and then show by induction that
$\psi_{j+1,\theta}$ is controlled by $\psi_{j,\theta}$ and $\delta_\theta$.
We use the constants of the threshold only at fixed $\theta$, before any
outer limit, and the argument does not require $L_{\mathrm{Lip},\theta}$ to stay
bounded.

\paragraph{Step 1: the initial span.}
Let $e_1,\dots,e_{k_0}$ be a $\zeta_N$-measurable orthonormal basis of
$W_0$, completed by zero vectors if $\dim W_0<k_0$. We obtain such a frame by Gram--Schmidt on the projected coordinate vectors,
keeping nonzero columns in a fixed order and then padding by zero columns. For $N\ge1$, fix a
deterministic unit vector $e_{\mathrm{ref}}\in\R^N$, and for every $i$ set
$w_i=\sqrt N e_i$ on $\{e_i\ne0\}$ and
$w_i=\sqrt N e_{\mathrm{ref}}$ on $\{e_i=0\}$. This gives a
$\zeta_N$-measurable direction of norm $\sqrt N$ on every seed event.
We apply \cref{sel:absolute-field-cost} with
$L_{\mathrm{Lip}}=L_{\mathrm{Lip},\theta}$ and direction $w_i$, and we drop
its nonnegative overlap penalty on $\{e_i=0\}$. Since the pairing with
$e_i$ vanishes there, this yields
\[
 \bd\limsup_{N\to\infty}\E\left|\frac{\ip{\mathcal A_{N,\theta}}{e_i}}{\sqrt N}\right|^p
 \le\Eth-\liminf_{N\to\infty}\E\frac{H_N(\mathcal A_{N,\theta})}N
 \le\delta_\theta.
\]
For each $i$, by Jensen's inequality and $p\ge2$, we have
$\E|\ip{\mathcal A_{N,\theta}}{e_i}|^2/N
\le(\E|\ip{\mathcal A_{N,\theta}}{e_i}/\sqrt N|^p)^{2/p}$. Summing over $i$,
we obtain
\begin{equation}\label{sel:adaptive-initial-estimate}
 \psi_{0,\theta}\le k_0\left(\frac{\delta_\theta}{\bd}\right)^{2/p}.
\end{equation}

\paragraph{Step 2: the energy of the residual.}
Fix $j<M$, write $W=W_j$, and let
$\mathcal A^\perp=P_{W^\perp}\mathcal A_{N,\theta}$. Both
$\mathcal A_{N,\theta}$ and $\mathcal A^\perp$ lie in $B_N$. Using \eqref{sel:tensor-ui} and the mean value theorem, we have
\[
 \frac{|H_N(\mathcal A_{N,\theta})-H_N(\mathcal A^\perp)|}N
 \le pQ_N\frac{\norm{P_W\mathcal A_{N,\theta}}}{\sqrt N}.
\]
By the Cauchy--Schwarz inequality, the expectation of the right side has
upper limit at most $pQ_*\sqrt{\psi_{j,\theta}}$. Combining this with the
definition \eqref{sel:adaptive-deficit} of $\delta_\theta$, we find
\begin{equation}\label{eq:residual-energy-lower}
 \liminf_{N\to\infty}\E\frac{H_N(\mathcal A^\perp)}N
 \ge\Eth-\delta_\theta-pQ_*\sqrt{\psi_{j,\theta}}.
\end{equation}

\paragraph{Step 3: the residual model.}
Since $k_0$ and $M$ are fixed before the dimension limit, we may restrict
to $N>k_0+M$. Every residual dimension below is then positive.
We now show that $\mathcal A^\perp$ is the output of a Lipschitz algorithm
for a pure $p$-spin model in dimension $N'=N-\dim W$, and that its
energy has the same normalization up to a factor tending to one. We
condition on $\mathcal F_j$. Then $W$ and $u_j$ are fixed, and we choose an
$\mathcal F_j$-measurable orthonormal basis $f_1,\dots,f_{N'}$ of
$W^\perp$. Let $J'\in(\R^{N'})^{\otimes p}$ be the coordinates of
$P_{\mathcal T(W)}J$ in the basis
$(f_{i_1}\otimes\dots\otimes f_{i_p})$. By
\cref{sel:adaptive-disclosure}\ref{item:disclosure-gaussian}, $J'$ is a
standard Gaussian vector in the $N'^p$-dimensional space
$(\R^{N'})^{\otimes p}$, conditionally on $\mathcal F_j$. Let
$y'\in\R^{N'}$ be the coordinates of $\mathcal A^\perp$ in the basis
$(f_i)$, and set $x=\sqrt{N'/N}\,y'$. Then $\norm x\le\sqrt{N'}$, that is,
$x\in B_{N'}$. Since $(\mathcal A^\perp)^{\otimes p}\in\mathcal T(W)$, only
$P_{\mathcal T(W)}J$ contributes to $H_N(\mathcal A^\perp)$, and a change of
normalization yields
\begin{equation}\label{sel:residual-normalization}
 \frac{H_N(\mathcal A^\perp)}N
 =N^{-(p+1)/2}\ip{J'}{y'^{\otimes p}}
 =\sqrt{\frac{N'}N}\,\frac{H'_{N'}(x)}{N'},
\end{equation}
where $H'_{N'}(x)=N'^{-(p-1)/2}\ip{J'}{x^{\otimes p}}$ is the pure $p$-spin
Hamiltonian in dimension $N'$ with coefficient tensor $J'$. Indeed,
$\ip{J'}{y'^{\otimes p}}=(N/N')^{p/2}\ip{J'}{x^{\otimes p}}$, and
$N^{-(p+1)/2}(N/N')^{p/2}N'^{(p-1)/2}=(N'/N)^{1/2}N'^{-1}$. Moreover, $J$
is an affine isometric function of $J'$ given $\mathcal F_j$, and
$\mathcal A^\perp\mapsto x$ is a contraction. Therefore, conditionally on $\mathcal F_j$, the map $J'\mapsto x$ is $L_{\mathrm{Lip},\theta}$-Lipschitz with
values in $B_{N'}$. Let $u'\in\R^{N'}$ be the coordinates of $u_j$ when
$u_j\ne0$, and the first coordinate vector otherwise. The direction
$w'=\sqrt{N'}u'$ has norm $\sqrt{N'}$, and when $u_j\ne0$,
\[
 \frac{\ip{w'}{x}}{N'}=\frac{\ip{u'}{y'}}{\sqrt N}
 =\frac{\ip{u_j}{\mathcal A_{N,\theta}}}{\sqrt N},
\]
where we used $u_j\in W^\perp$ in the last equality.

\paragraph{Step 4: the penalty in the residual model.}
Let $\alpha_j=\ip{u_j}{\mathcal A_{N,\theta}}/\sqrt N$, which is zero if
$u_j=0$. The conditional law of $J'$ given $\mathcal F_j$ is standard
Gaussian, and on each event $\{N'=n'\}$ the tensor $J'$ is independent of
$\mathcal F_j$. We may therefore apply \eqref{eq:conditional-field-cost} in
dimension $N'$, with $\mathcal F_j$ in the role of the seed, to the
algorithm $J'\mapsto x$ and the direction $w'$. If $u_j\ne0$, its overlap
term is $\bd|\alpha_j|^p$; if $u_j=0$, we drop its nonnegative overlap
term. Since $N'\ge N-k_0-j$ and $\phi_N\ge0$, we obtain
\[
 \E\left[\frac{H'_{N'}(x)}{N'}+\bd|\alpha_j|^p\,\middle|\,\mathcal F_j\right]
 \le\Eth+\sup_{N''\ge N-k_0-j}\phi_{N''}(p,L_{\mathrm{Lip},\theta}).
\]
We use \eqref{sel:residual-normalization} to write
\[
 \frac{H_N(\mathcal A^\perp)}N+\bd|\alpha_j|^p
 =\sqrt{\frac{N'}N}\left(\frac{H'_{N'}(x)}{N'}+\bd|\alpha_j|^p\right)
 +\left(1-\sqrt{\frac{N'}N}\right)\bd|\alpha_j|^p.
\]
Here $\sqrt{N'/N}\le1$, $0\le1-\sqrt{N'/N}\le(k_0+j)/N$, and $|\alpha_j|\le1$.
Taking expectations and letting $N\to\infty$, we find
\[
 \limsup_{N\to\infty}\left\{\E\frac{H_N(\mathcal A^\perp)}N
 +\bd\,\E|\alpha_j|^p\right\}\le\Eth.
\]
Combining this with \eqref{eq:residual-energy-lower}, we obtain
$\bd\limsup_N\E|\alpha_j|^p\le\delta_\theta+pQ_*\sqrt{\psi_{j,\theta}}$. By
Jensen's inequality as in Step~1, and since
$\norm{P_{W_{j+1}}\mathcal A_{N,\theta}}^2=\norm{P_{W_j}\mathcal A_{N,\theta}}^2
+N|\alpha_j|^2$, we conclude that
\begin{equation}\label{sel:adaptive-recursion}
 \psi_{j+1,\theta}\le\psi_{j,\theta}
 +\left(\frac{\delta_\theta+pQ_*\sqrt{\psi_{j,\theta}}}{\bd}\right)^{2/p}.
\end{equation}

\paragraph{Step 5: the outer limits.}
The right sides of \eqref{sel:adaptive-initial-estimate} and
\eqref{sel:adaptive-recursion} are continuous nondecreasing functions of
$(\delta_\theta,\psi_{j,\theta})$ that vanish at $(0,0)$. Upper limits pass
through such functions. Taking the outer limits in their prescribed order
and inducting on $j\le M$, we find
$\lim_\theta\psi_{j,\theta}=0$ for every $j\le M$. This is
\eqref{sel:finite-span-orthogonality}. Finally, if
$Y_N\in W_M\cap B_N$, then
\[
 \left|\frac{\ip{\mathcal A_{N,\theta}}{Y_N}}N\right|^2
 =\left|\frac{\ip{P_{W_M}\mathcal A_{N,\theta}}{Y_N}}N\right|^2
 \le\frac{\norm{P_{W_M}\mathcal A_{N,\theta}}^2}N,
\]
which proves the last assertion.
\end{proof}

\section{Transfer to a fixed waiting-time state}\label{sec:general-transfer}

We now combine the forks with finite-query orthogonality on their fixed
product probability spaces. The hypotheses of the next proposition are those of
\cref{prop:threshold-inputs}\ref{item:threshold-forgetting}. Under them,
the translated family $(C_T,R_T,\lambda_T)$ converges along every sequence
$T_n\to\infty$ to one stationary cluster $(c,r,\ell)$. Indeed, by the
precompactness in \cref{prop:stationary-clusters}, every subsequence of
$(T_n)$ has a further subsequence along which the family converges, and the
limit is a stationary cluster whose kernels are $(c,r)$. Its multiplier is
determined by $r$: setting $\tau=0$ in \eqref{eq:stationary-cluster-response}
and using $r(0)=1$, we find $\ell=-r'(0)$. All subsequential limits are
therefore equal. In particular, \cref{sel:finite-forks} applies along every
sequence $T_n\to\infty$.

\begin{proposition}[Fixed-waiting-time decorrelation]
\label{sel:general-fixed-wait}
Fix $p\ge3$ and $\beta>0$. Suppose that the translated kernels converge
along every sequence $T\to\infty$ to one stationary profile $(c,r)$, and
that its plateau $q$ is positive and satisfies
$\beta^2p(p-1)q^{p-2}(1-q)^2=1$. Then
\[
 \lim_{t\to\infty}C(t,s)=0\qquad\text{for every fixed }s\ge0.
\]
\end{proposition}

\begin{proof}
The proof is by contradiction. Recall from \eqref{eq:ck-physical-bounds}
that $C\ge0$. Suppose that, for some $s\ge0$, there exist
$\gamma>0$ and $t_n\to\infty$ such that $C(t_n,s)\ge\gamma$ for all $n$. The
proof has three steps. We first express $C(t_n,s)$ through the center of a
fork that ends at $t_n$. Then we show that the Lipschitz algorithms
approximating the normalized center saturate the threshold in the order
\eqref{sel:aging-limit-order}. Finally, on the fixed fork product law, we approximate $X_s$ by a
finite-query Euler vector and use \cref{sel:adaptive-orthogonality} to force
their overlap with $X_s$ to vanish, which contradicts the first step.

\paragraph{Step 1: the overlap of the center.}
Fix $K\ge1$ and $L\ge0$, and set $T_n=\max\{s+1,t_n-L\}$ for every $n$.
Then $T_n\to\infty$, $T_n\ge s$, and for all sufficiently large $n$ we have
$T_n+L=t_n$. Let $m_{N,n}$ be the center of the fork at time $T_n$ with $K$
branches of length $L$. For these sufficiently large $n$, the branches
end at time $t_n$. The time $s$ lies in the common past,
and each pair $(Y_b,X_s)$ has the law of $(X_{t_n},X_s)$. This implies that
\begin{equation}\label{sel:fork-fixed-wait-overlap}
 \E\frac{\ip{m_{N,n}}{X_s}}N=\E C_N(t_n,s),
\end{equation}
and by \cref{prop:hard-bridge}, the right side tends to $C(t_n,s)$ as
$N\to\infty$. In the limit in $n$, the endpoint of the fork is the time $t_n$ itself,
so we do not need any regularity of $C(\cdot,s)$ in its first variable.

\paragraph{Step 2: saturation of the threshold.}
Let $\theta=(n,\varepsilon,K,L)$, where $\varepsilon=1/(m+1)$ and
$K=k+1$ for $m,k\in\mathbb N$, with the order of
\eqref{sel:aging-limit-order}, and let
$\mathcal A_{N,\theta}=\mathcal A^\varepsilon_{N,T_n}$ be the algorithms of
\cref{sel:finite-forks}\ref{item:fork-lipschitz} for this fork. They satisfy
the standing hypotheses of \cref{sec:general-disclosure} on the fixed fork
product law of \cref{sec:general-forks}, with seed $\zeta_N$ consisting
of $X_0$, the $K$ full future Wiener paths, and the full trunk Wiener path $B$,
jointly independent of $J$, and with a constant $L_{\mathrm{Lip},\theta}$
depending only on $p$, $\beta$, the fixed waiting time $s$, the horizon
$T_n+L$, $K$, and $\varepsilon$. We write
$Q=Q_{K,L}$, and recall that $q\le Q\le1$. Let
\[
 \Delta_{N,\theta}=\left(\frac1N\E\norm{\mathcal A_{N,\theta}
 -m_{N,n}/\sqrt Q}^2\right)^{1/2}.
\]
By \eqref{sel:fork-lipschitz} and the comparison of normalizations in
\cref{sel:finite-forks}\ref{item:fork-lipschitz}, we have
\begin{equation}\label{eq:transfer-approximation}
 \limsup_{n\to\infty}\limsup_{j\to\infty}\Delta_{N_j,\theta}\le\varepsilon.
\end{equation}
The vector $m_{N,n}/\sqrt Q$ has norm at most $\sqrt{N/q}$. Using
\eqref{eq:general-energy-transfer} with $R=q^{-1/2}$ and the homogeneity of
$H_N$, we find
\[
 \E\frac{H_N(\mathcal A_{N,\theta})}N
 \ge Q^{-p/2}\E\frac{H_N(m_{N,n})}N-C_q\Delta_{N,\theta},\qquad
 C_q=p\,q^{-(p-1)/2}Q_*.
\]
The constant $C_q$ does not depend on $n$, $\varepsilon$, $K$, or $L$; this
is where we use $q>0$. Then \eqref{sel:fork-center-energy-limit} implies that
$\lim_n\liminf_jQ^{-p/2}\E H_{N_j}(m_{N_j,n})/N_j=Q^{-p/2}h_{K,L}$.
Together with \eqref{eq:transfer-approximation}, this gives
\[
 \limsup_{n\to\infty}\delta_\theta
 \le\max\left\{0,\;\Eth-\frac{h_{K,L}}{Q_{K,L}^{p/2}}+C_q\varepsilon\right\}.
\]
We let $\varepsilon\downarrow0$, then $K\to\infty$, then $L\to\infty$. By
\eqref{sel:fork-saturation}, which applies because $q$ is marginal, we
conclude that $\lim_\theta\delta_\theta=0$ in the order
\eqref{sel:aging-limit-order}.

\paragraph{Step 3: orthogonality on the fork product law and the contradiction.}
We use the algorithms of Step~2 on precisely the fork seed and disorder
product space just described, and take the limits in precisely the order
\eqref{sel:aging-limit-order}, with $\varepsilon=1/(m+1)$. All disclosure
histories below include the full fork seed. If $s=0$, then $X_0$ lies in the
seed-measurable subspace $W_0=\Span(X_0)$, and
\cref{sel:adaptive-orthogonality} with $k_0=1$ and no disclosure steps gives
zero limiting squared overlap. For $s>0$, fix $M\ge1$ before any limit.
We construct the Euler--Maruyama approximation
\cite[Sections~4--5]{higham2001} of this fork's trunk state $X_s$ and prove
the required normalized error directly.

\subparagraph{The coefficients.}
Let $\Psi_N=\Pi_{B_N}$ be the metric projection onto the ball of radius $\sqrt N$.
We use $P_y=I-yy^{\mathsf T}/N$ for every $y\in\R^N$, as in \cref{sec:model}.
For $x\in\R^N$, we define
\[
 b_N(x)=\beta P_{\Psi_N(x)}\nabla H_N(\Psi_N(x))-\frac{N-1}N\Psi_N(x),\qquad
 \sigma_N(x)=\sqrt2\,P_{\Psi_N(x)}.
\]
Since $\Psi_N(x)=x$ on $S_N$, these are the coefficients of
\eqref{eq:hard-sde} on $S_N$, and
\[
 X_t=X_0+\int_0^tb_N(X_r)\,\dd r+\int_0^t\sigma_N(X_r)\,\dd B_r.
\]
Let $G_N=\{Q_N\le A_p\}$, and let $y=\Psi_N(x)$ and $z=\Psi_N(x')$, so that
$\norm y,\norm z\le\sqrt N$ and $\norm{y-z}\le\norm{x-x'}$. The matrix $P_y$
has eigenvalues $1$ and $1-\norm y^2/N\in[0,1]$, which gives
$\norm{P_y}_\op\le1$ and $\norm{P_y}_{\HS}^2\le N$. Moreover,
\[
 \norm{P_y-P_z}_{\HS}=\frac{\norm{yy^{\mathsf T}-zz^{\mathsf T}}_{\HS}}N
 \le\frac{(\norm y+\norm z)\norm{y-z}}N\le\frac{2\norm{x-x'}}{\sqrt N}.
\]
On $G_N$, \eqref{eq:general-ball-derivatives} with $R=2$ gives
$\norm{\nabla H_N(y)}\le p2^{p-1}A_p\sqrt N$ and
$\norm{\nabla H_N(y)-\nabla H_N(z)}\le p(p-1)2^{p-2}A_p\norm{y-z}$. Writing
$P_y\nabla H_N(y)-P_z\nabla H_N(z)
=P_y(\nabla H_N(y)-\nabla H_N(z))+(P_y-P_z)\nabla H_N(z)$, we obtain, on
$G_N$ and for all $x,x'\in\R^N$,
\begin{equation}\label{eq:euler-coefficient-bounds}
\begin{gathered}
 \norm{b_N(x)-b_N(x')}\le K_b\norm{x-x'},\qquad
 \norm{\sigma_N(x)-\sigma_N(x')}_{\HS}\le\frac{2\sqrt2}{\sqrt N}\norm{x-x'},\\
 \norm{b_N(x)}^2\le K_b'N,\qquad\norm{\sigma_N(x)}_{\HS}^2\le2N,
\end{gathered}
\end{equation}
with $K_b=1+\beta pA_p2^{p-2}(3p+5)$ and
$K_b'=(3\beta p2^{p-1}A_p+2)^2$. The Hilbert--Schmidt bound on
$\sigma_N(x)-\sigma_N(x')$ is the reason that the error below does not grow with $N$: a bound of order one for each matrix entry would lose a factor of
$N$.

\subparagraph{The Euler--Maruyama scheme.}
Let $t_i=is/M$ for $0\le i\le M$. We define $Z_0=X_0$ and
\[
 Z_{i+1}=Z_i+\frac sM\,b_N(Z_i)+\sigma_N(Z_i)\big(B_{t_{i+1}}-B_{t_i}\big),
 \qquad 0\le i<M.
\]
For $r\in[t_i,t_{i+1})$, let $r_-=t_i$, and let
$Z_r=Z_i+(r-t_i)b_N(Z_i)+\sigma_N(Z_i)(B_r-B_{t_i})$, so that
\[
 Z_t=X_0+\int_0^tb_N(Z_{r_-})\,\dd r+\int_0^t\sigma_N(Z_{r_-})\,\dd B_r
\]
and $Z_s=Z_M$. We set $Y_{N,M}=\Pi_{B_N}(Z_M)$.

\subparagraph{The error on $G_N$.}
Fix $J\in G_N$ and $X_0$, and let $\E_B$ denote expectation over the
Brownian motion. Let $u(t)=\E_B\norm{X_t-Z_t}^2/N$ and
$v(r)=\E_B\norm{Z_r-Z_{r_-}}^2/N$. Since $Z_{r_-}$ is independent of
$B_r-B_{r_-}$, the It\^o isometry and \eqref{eq:euler-coefficient-bounds}
give
\[
 v(r)\le2\Big(\frac sM\Big)^2K_b'+2\,\frac sM\,\frac{2N}N
 \le C_1\frac sM,\qquad C_1=2sK_b'+4.
\]
Subtracting the integral equations for $X_t$ and $Z_t$, and using the
Cauchy--Schwarz inequality for the drift and the It\^o isometry for the
noise, we find, for $t\le s$,
\[
 u(t)\le2sK_b^2\int_0^t\E_B\frac{\norm{X_r-Z_{r_-}}^2}N\,\dd r
 +\frac{16}N\int_0^t\E_B\frac{\norm{X_r-Z_{r_-}}^2}N\,\dd r.
\]
Inserting the bound $\E_B\norm{X_r-Z_{r_-}}^2/N\le2u(r)+2v(r)$, we obtain
$u(t)\le C_2\int_0^tu(r)\,\dd r+C_2C_1s^2/M$ with $C_2=4sK_b^2+32$. By
Gronwall's inequality,
\[
 u(s)\le C_1C_2s\,e^{C_2s}\,\frac sM.
\]
The projection $\Pi_{B_N}$ is $1$-Lipschitz and fixes $X_s\in S_N$, which
gives
$\E_B\norm{X_s-Y_{N,M}}^2/N\le u(s)$ on $G_N$. On $G_N^c$, we have
$\norm{X_s-Y_{N,M}}^2/N\le4$, and $\Prob(G_N^c)\le e^{-c_pN}$. Thus, with $C=C_1C_2se^{C_2s}$,
\begin{equation}\label{sel:euler-approximation}
 \limsup_{N\to\infty}\E\frac{\norm{X_s-Y_{N,M}}^2}N\le C\,\frac sM.
\end{equation}
The expectation here is on the fixed fork product law. Its projection
onto $(J,X_0,B)$ has exactly the original disorder, uniform initial and
Wiener product law, so the error bound is independent of $n$, $\varepsilon$,
$K$, and $L$.

\subparagraph{The query span.}
Let $W_0$ be the span of $X_0$ and the increments $B_{t_{i+1}}-B_{t_i}$,
$0\le i<M$. It is measurable with respect to the full fork seed $\zeta_N$ and has dimension at
most $M+1$. Suppose that $Z_i\in W_i$ and that $Z_i$ is
$\mathcal F_i$-measurable, where $\mathcal F_i$ is the disclosure history
generated using this full fork seed, which holds for $i=0$. The point $\Psi_N(Z_i)$
is a multiple of $Z_i$ and lies in $W_i$. By \cref{sel:adaptive-disclosure}\ref{item:disclosure-gradient}, the gradient
$g_i=\nabla H_N(\Psi_N(Z_i))$ is $\mathcal F_i$-measurable. We choose $u_i$
as the normalized projection of $g_i$ onto $W_i^\perp$ (zero if it
vanishes), so that $g_i\in W_{i+1}$. Since
$P_yw=w-y\ip yw/N$, the step from $Z_i$ to $Z_{i+1}$ is a linear
combination of $Z_i$, $\Psi_N(Z_i)$, $g_i$, and $B_{t_{i+1}}-B_{t_i}$, with
coefficients that are $\mathcal F_i$-measurable. This shows that
$Z_{i+1}\in W_{i+1}$ and that $Z_{i+1}$ is $\mathcal F_{i+1}$-measurable. After $M$ steps, $Z_M\in W_M$,
and $Y_{N,M}$, a multiple of $Z_M$, lies in $W_M$. Each step adds at most
one dimension, which gives $\dim W_M\le2M+1$.

\subparagraph{Passing from the Euler vector to the waiting state.}
By Step~2, \cref{sel:adaptive-orthogonality} applies to these particular
full-fork disclosure histories, with $k_0=M+1$ and $M$ queries. Hence the
squared overlap of $\mathcal A_{N,\theta}$ with $Y_{N,M}$ vanishes after the
dimension limit and then the limits in $n$, $\varepsilon=1/(m+1)$, $K$, and $L$.
Since $\norm{\mathcal A_{N,\theta}}\le\sqrt N$, we have
\[
 \E\left|\frac{\ip{\mathcal A_{N,\theta}}{X_s}}N\right|^2
 \le2\E\left|\frac{\ip{\mathcal A_{N,\theta}}{Y_{N,M}}}N\right|^2
 +2\E\frac{\norm{X_s-Y_{N,M}}^2}N.
\]
By \eqref{sel:euler-approximation}, the nested upper limit of the left side
is at most $2Cs/M$. It does not depend on $M$. Sending $M\to\infty$ last
(and using the $s=0$ argument above when needed), we obtain
\begin{equation}\label{eq:transfer-orthogonality}
 \lim_\theta\limsup_{N\to\infty}
 \E\left|\frac{\ip{\mathcal A_{N,\theta}}{X_s}}N\right|^2=0.
\end{equation}
On the other hand, by \eqref{eq:fork-overlap-transfer} with $V_N=X_s$ and by
\eqref{sel:fork-fixed-wait-overlap}, we have
\[
 \frac{\E C_N(t_n,s)}{\sqrt Q}
 \le\left(\E\left|\frac{\ip{\mathcal A_{N,\theta}}{X_s}}N\right|^2\right)^{1/2}
 +\Delta_{N,\theta}.
\]
As $N=N_j\to\infty$, the left side tends to $C(t_n,s)/\sqrt Q$, which is at
least $\gamma$ because $Q\le1$. We conclude that, for every fixed
$\varepsilon$, $K$, and $L$ and all sufficiently large $n$,
\[
 \gamma\le\limsup_{j\to\infty}
 \left(\E\left|\frac{\ip{\mathcal A_{N_j,\theta}}{X_s}}{N_j}\right|^2\right)^{1/2}
 +\limsup_{j\to\infty}\Delta_{N_j,\theta}.
\]
We take the upper limit as $n\to\infty$, then let $\varepsilon\downarrow0$, $K\to\infty$, and $L\to\infty$, in the order \eqref{sel:aging-limit-order}.
By \eqref{eq:transfer-orthogonality}, the first term tends to zero, and by
\eqref{eq:transfer-approximation}, the second is at most $\varepsilon$ after
the limit in $n$. Hence $\gamma\le0$, which contradicts $\gamma>0$. This completes the proof.
\end{proof}

\section{Proof of \texorpdfstring{\cref{prop:threshold-inputs}}{the threshold inputs}}
\label{sec:general-proof-inputs}

We now prove both parts of \cref{prop:threshold-inputs} for every $p\ge3$.

\begin{proof}[Proof of \cref{prop:threshold-inputs}\ref{item:threshold-ceiling}]
Let $T_n\to\infty$ be a sequence along which the translated kernels converge
to a stationary cluster $(c,r,\ell)$ with plateau $q>0$, and fix $L\ge0$.
For every $K\ge1$, \cref{sel:finite-forks} applies along $T_n$, and
\eqref{sel:fork-center-threshold} gives
\[
 h_{K,L}\le\Eth\,Q_{K,L}^{p/2}.
\]
By \cref{lem:fork-parameter-limits}\ref{item:fork-K-limit}, we have
$|Q_{K,L}-d_L|\le1/K$ and $|h_{K,L}-h_L|\le\beta p(p-1)/K$. The second bound
is the dominated convergence argument made explicit: the difference
$h_{K,L}-h_L$ is $\beta p$ times the integral of $r(\tau)$ times
$k_{K,L}(\tau)^{p-1}-c(2L-\tau)^{p-1}$, which is at most $(p-1)/K$ in absolute
value, and $\int_0^Lr\le1$. Letting $K\to\infty$ and using the continuity of
$x\mapsto x^{p/2}$, we obtain
\[
 h_L\le\Eth\,d_L^{p/2},
\]
which is \eqref{eq:threshold-input-ceiling}.
\end{proof}

Letting $L\to\infty$ in \eqref{eq:threshold-input-ceiling} and using
\cref{lem:fork-parameter-limits}\ref{item:fork-L-limit}, we obtain
\begin{equation}\label{sel:general-positive-upper-bound}
 h\le\Eth\,q^{p/2},
\end{equation}
where $h$ is given by \eqref{sel:limiting-center-energy}. This is the form
in which part~\ref{item:threshold-ceiling} enters the proof of
\cref{sel:positive-marginality}. It holds for every stationary cluster with
positive plateau, and its proof does not use marginality.

\begin{proof}[Proof of \cref{prop:threshold-inputs}\ref{item:threshold-forgetting}]
The hypotheses of part~\ref{item:threshold-forgetting} are those of
\cref{sel:general-fixed-wait}, which gives $C(t,s)\to0$ as $t\to\infty$ for
each fixed $s$. It remains to make the convergence uniform in
$s\in[0,S]$. Let $k\ge1$ be an integer and let $s_i=i/k$ for
$0\le i\le\lceil kS\rceil$. Every $s\in[0,S]$ is within $1/k$ of some
$s_i$, and by the correlation modulus \eqref{eq:global-modulus} in the
second variable,
\[
 |C(t,s)-C(t,s_i)|\le\frac{D_{p,\beta}}k+\sqrt{\frac2k}
\]
for all $t$. Since $C\ge0$, we obtain
\[
 \sup_{0\le s\le S}|C(t,s)|\le\max_{i}C(t,s_i)+\frac{D_{p,\beta}}k+\sqrt{\frac2k}.
\]
The maximum is over finitely many points, and it tends to zero as
$t\to\infty$ by \cref{sel:general-fixed-wait}. Letting $t\to\infty$ and then
$k\to\infty$ proves \eqref{eq:threshold-input-forgetting}.
\end{proof}
